\subsection{Extension of Scalars in Semisimple Modules}

PROPOSITION 8. --- Let A be a K-algebra, M an A-module, and L an extension of the field K. Denote the commutant of M by D and the center of D by Z.

\begin{enumerate}
\def\labelenumi{\alph{enumi})}
\item
  Suppose that the \(A_{(L)}\) -module \(M_{(L)}\) is simple (resp. isotypical, semisimple). Then the A-module M is simple (resp. isotypical, semisimple).
\item
  Suppose that the A-module M is semisimple and that M or L is finite-dimensional over K. The \(\mathrm{A}_{(\mathrm{L})}\) -module \(\mathrm{M}_{(\mathrm{L})}\) is semisimple if and only if the ring \(\mathrm{Z}_{(\mathrm{L})}\) is reduced. The \(\mathrm{A}_{(\mathrm{L})}\) -module \(\mathrm{M}_{(\mathrm{L})}\) is isotypical and nonzero if and only if the ring \(\mathrm{Z}_{(\mathrm{L})}\) is an integral domain.
\item
  Suppose that the A-module M is simple. The \(\mathrm{A}_{(\mathrm{L})}\) -module \(\mathrm{M}_{(\mathrm{L})}\) is semisimple (resp. isotypical, simple) if and only if the ring \(\mathrm{D}_{(\mathrm{L})}\) is semisimple (resp. simple, a field).
\end{enumerate}

Assertion a) is a specific case of Proposition 1 (VIII, p.~211), assertion b) is a specific case of Proposition 5 (VIII, p.~219), and assertion c) is a specific case of Corollary 1 of VIII, p.~215.

COROLLARY 1. --- a) Suppose that the A-module M is semisimple, that the extension L of K is separable, and that M or L is finite-dimensional over K. Then the \(\mathrm{A}_{(\mathrm{L})}\) -module \(\mathrm{M}_{(\mathrm{L})}\) is semisimple.

\begin{enumerate}
\def\labelenumi{\alph{enumi})}
\setcounter{enumi}{1}
\tightlist
\item
  Suppose that the A-module M is simple and that its commutant is equal to K. Then the \(\mathrm{A}_{(\mathrm{L})}\) -module \(\mathrm{M}_{(\mathrm{L})}\) is simple.
\end{enumerate}

Assertion a) follows from Corollary VIII, p.~221. Assertion b) is a specific case of Proposition 8, c).

COROLLARY 2. --- Let L be an extension of the field K. Denote the center of the K-algebra A by Z.

\begin{enumerate}
\def\labelenumi{\alph{enumi})}
\item
  If the L-algebra \(A_{(L)}\) is semisimple, then the K-algebra A is semisimple.
\item
  Suppose that the K-algebra A is semisimple and that L or A is finite-dimensional over K. The L-algebra \(\mathrm{A}_{(\mathrm{L})}\) is semisimple if and only if the ring \(\mathrm{Z}_{(\mathrm{L})}\) is reduced; this is, in particular, the case when L is a separable extension of K. The L-algebra \(\mathrm{A}_{(\mathrm{L})}\) is simple if and only if the ring \(\mathrm{Z}_{(\mathrm{L})}\) is an integral domain; this is, in particular, the case when the center of A is equal to K.
\end{enumerate}

Assertions a) and b) follow from Proposition 8, a) and b) applied to the A-module \(\mathrm{A}_s\) .

PROPOSITION 9. --- Let A be a K-algebra and L a separable extension of K.

\begin{enumerate}
\def\labelenumi{\alph{enumi})}
\item
  If M is an A-module without radical, then the \(\mathrm{A}_{(\mathrm{L})}\) -module \(\mathrm{M}_{(\mathrm{L})}\) is without radical.
\item
  If the K-algebra A is without radical, then the L-algebra \(\mathrm{A}_{(\mathrm{L})}\) is without radical.
\end{enumerate}

Let us prove assertion a). Let M be an A-module without radical. We identify M with its canonical image in \(\mathrm{M}_{(\mathrm{L})}\) . Let N be a maximal submodule of M. Since the A-module M/N is simple, it follows from Proposition 4 of VIII, p.~218 that the \(\mathrm{A}_{(\mathrm{L})}\) -module \((\mathrm{M}/\mathrm{N})_{(\mathrm{L})} = \mathrm{M}_{(\mathrm{L})}/\mathrm{N}_{(\mathrm{L})}\) is without radical and therefore \(\mathfrak{R}_{\mathrm{A}_{(\mathrm{L})}}(\mathrm{M}_{(\mathrm{L})}) \subset \mathrm{N}_{(\mathrm{L})}\) by Corollary 1, c) of VIII, p.~152. Now, it follows from the corollary of Proposition 14 of II, §7, No.~7, p.~306 that the intersection of the \(\mathrm{N}_{(\mathrm{L})}\) , where N runs through the set of maximal submodules of M, is reduced to 0. Consequently, the \(\mathrm{A}_{(\mathrm{L})}\) -module \(\mathrm{M}_{(\mathrm{L})}\) is without radical.

Assertion b) follows from assertion a) applied to the A-module \(\mathrm{A}_s\) .

PROPOSITION 10. --- Let A be a K-algebra and L an extension of K. Let M be an A-module.

\begin{enumerate}
\def\labelenumi{\alph{enumi})}
\tightlist
\item
  We make one of the following two assumptions:
\end{enumerate}

\begin{enumerate}
\def\labelenumi{(\roman{enumi})}
\item
  The A-module M is finitely generated and L is algebraic over K.
\item
  The ring A is left Artinian.
\end{enumerate}

Then we have the inclusion

\[
\mathfrak {R} _ {\mathrm{A}} (\mathrm{M}) _ {(\mathrm{L})} \subset \mathfrak {R} _ {\mathrm{A} _ {(\mathrm{L})}} (\mathrm{M} _ {(\mathrm{L})}).
\]

\begin{enumerate}
\def\labelenumi{\alph{enumi})}
\setcounter{enumi}{1}
\tightlist
\item
  If L is a separable extension of K, then we have the inclusion
\end{enumerate}

\[
\mathfrak {R} _ {\mathrm{A} _ {(\mathrm{L})}} (\mathrm{M} _ {(\mathrm{L})}) \subset \mathfrak {R} _ {\mathrm{A}} (\mathrm{M}) _ {(\mathrm{L})}.
\]

Let us prove assertion a). Consider case (i). First suppose that L has finite degree over K. Then the A-module \(M_{(L)}\) is finitely generated. Denote the canonical homomorphism from A to \(A_{(L)}\) by f; the ring \(A_{(L)}\) is generated by the union of its center and \(f(A)\) . We can therefore apply Proposition 3 of VIII, p.~174 to the \(A_{(L)}\) -module \(M_{(L)}\) . This gives the inclusion \(\mathfrak{R}_{\mathrm{A}}(\mathrm{M}_{(\mathrm{L})}) \subset \mathfrak{R}_{\mathrm{A}_{(\mathrm{L})}}(\mathrm{M}_{(\mathrm{L})})\) and a fortiori \(\mathfrak{R}_{\mathrm{A}}(\mathrm{M})_{(\mathrm{L})} \subset \mathfrak{R}_{\mathrm{A}_{(\mathrm{L})}}(\mathrm{M}_{(\mathrm{L})})\) (VIII, p.~152, Corollary 1).

Let us now treat the general case. Let \(x_{1},\ldots,x_{n}\) be elements that generate the A-module M and x an element of \(\mathfrak{R}_{\mathrm{A}}(\mathrm{M})\) . Let \(a_{1},\ldots,a_{n}\) be elements of \(\mathrm{A}_{(\mathrm{L})}\) ; since L is algebraic over K, there exists a finite extension \(L'\) of K contained in L such that the \(a_{i}\) belong to \(\mathrm{A}_{(\mathrm{L}')}\) . By the above, x belongs to the radical of the \(\mathrm{A}_{(\mathrm{L}')}\) -module \(\mathrm{M}_{(\mathrm{L}')}\) ; it follows from the corollary of VIII, p.~153 that the elements \(x_{i}+a_{i}x\) for \(1\leqslant i\leqslant n\) generate the \(\mathrm{A}_{(\mathrm{L}')}\) -module \(\mathrm{M}_{(\mathrm{L}')}\) and therefore the \(\mathrm{A}_{(\mathrm{L})}\) -module \(\mathrm{M}_{(\mathrm{L})}\) . By the same corollary, x belongs to the radical of the \(\mathrm{A}_{(\mathrm{L})}\) -module \(\mathrm{M}_{(\mathrm{L})}\) .

Let us now consider case (ii). Let \(\mathfrak{r}\) be the radical of A, so that the radical of the A-module M is equal to \(\mathfrak{rM}\) (VIII, p.~174, Corollary). The radical \(\mathfrak{r}\) of A is a nilpotent two-sided ideal of A (VIII, p.~173, Proposition 1); therefore, \(\mathfrak{r}_{(L)}\) is a nilpotent two-sided ideal of \(A_{(L)}\) . It follows that \(\mathfrak{r}_{(L)} \subset \Re(A_{(L)})\) (VIII, p.~156, Theorem 1), and Proposition 6 of VIII, p.~158 implies \(\Re(A_{(L)})M_{(L)} \subset \Re_{A_{(L)}}(M_{(L)})\) . We therefore have \(\Re_A(M) = \mathfrak{rM} \subset \Re_{A_{(L)}}(M_{(L)})\) , which concludes the proof of assertion a).

The module \(\mathrm{M}/\mathfrak{R}_{\mathrm{A}}(\mathrm{M})\) is without radical. If L is a separable extension of K, then it follows from Proposition 9 of VIII, p.~223 that the \(\mathrm{A}_{(\mathrm{L})}\) -module \((\mathrm{M}/\mathfrak{R}_{\mathrm{A}}(\mathrm{M}))_{(\mathrm{L})}\) is without radical. Consequently, we have the inclusion

\[
\mathfrak {R} _ {\mathrm{A} _ {(\mathrm{L})}} (\mathrm{M} _ {(\mathrm{L})}) \subset \mathfrak {R} _ {\mathrm{A}} (\mathrm{M}) _ {(\mathrm{L})}.
\]

COROLLARY. --- Let L be a separable extension of K. We have \(\mathfrak{R}(A_{(L)}) = \mathfrak{R}(A)_{(L)}\) if L is algebraic over K or if the ring A is left Artinian.

This is the case \(\mathrm{M} = \mathrm{A}_s\) of Proposition 10.

\BourbakiExercises{Exercises}

\begin{enumerate}
\def\labelenumi{\arabic{enumi})}
\tightlist
\item
  Let A and B be K-algebras and M an A-module; view \(M \otimes B\) as a module over the algebra \(C = A \otimes B\) . Identify M with an A-submodule of \(M \otimes B\) .
\end{enumerate}

\begin{enumerate}
\def\labelenumi{\alph{enumi})}
\item
  Prove the inclusion \(M \cap \mathfrak{R}_{\mathrm{C}}(\mathrm{M} \otimes \mathrm{B}) \subset \mathfrak{R}_{\mathrm{A}}(\mathrm{M})\) (consider an A-linear homomorphism from M to a simple A-module).
\item
  Prove that we have equality in each of the following cases:
\end{enumerate}

\begin{enumerate}
\def\labelenumi{(\roman{enumi})}
\item
  A is finite-dimensional over \(\mathbf{K}\) .
\item
  The A-module M is finitely generated, and A is the union of a right directed family of subalgebras that are finite-dimensional over K.
\item
  \(M = A_{s}\) , and the radical of A is nilpotent.
\end{enumerate}

(In case (i), one proves that every homomorphism from \(\mathrm{M} \otimes \mathrm{B}\) to a simple C-module S is zero on \(\Re_{\mathrm{A}}(\mathrm{M})\) , by observing that \(\Re_{\mathrm{C}}(\mathrm{S})\) is zero.)

\begin{enumerate}
\def\labelenumi{\arabic{enumi})}
\setcounter{enumi}{1}
\item
  Let \(\mathrm{A}\) be a K-algebra that is an integral domain and \(\mathrm{L}\) its field of fractions. Prove that \(\mathrm{L}\) is a simple \(\mathrm{A} \otimes \mathrm{L}\) -module but that its radical as an A-module is equal to \(\mathrm{L}\) .
\item
  Give an example of (nonzero) K-algebras \(\mathrm{A}_1\) and \(\mathrm{A}_2\) such that \(\Re (\mathrm{A}_1)\neq 0\) and \(\Re (\mathrm{A}_1\otimes \mathrm{A}_2) = 0\) (cf.~VIII, p.~167, Exercise 11).
\item
  Give an example of a commutative field K and two extensions E and F of K such that the ring \(A = E \otimes_{K} F\) is semisimple but is not a field (resp. is not semisimple). Deduce that the A-module \(A_{s} = E_{s} \otimes_{K} F_{s}\) is semisimple but not simple (resp. is not semisimple).
\end{enumerate}

¶ 5) Let A be a quasi-simple K-algebra (VIII, p.~120, Remark 2).

\begin{enumerate}
\def\labelenumi{\alph{enumi})}
\item
  Prove that the (A, A)-bimodule \(_{s}A_{d}\) is simple; deduce that the center Z of A is a field, isomorphic to the commutant of this bimodule.
\item
  Let B be a K-algebra. Prove that every two-sided ideal of \(A \otimes B\) is of the form \(A \otimes_{Z} b\) , where b is a two-sided ideal of B (reduce to the case Z = K, and use Corollary 2 of VIII, p.~63).
\item
  Deduce that if A and B are quasi-simple K-algebras of which one has center K, then the algebra \(\mathrm{A} \otimes \mathrm{B}\) is quasi-simple.
\end{enumerate}

\(d\) ) Let B be a K-algebra and A a quasi-simple subalgebra of B whose center Z contains K. Prove that A and its commutant \(\mathrm{A}'\) in B are linearly disjoint over Z (use \(b\) ).

\begin{enumerate}
\def\labelenumi{\alph{enumi})}
\setcounter{enumi}{4}
\item
  Let A be a quasi-simple K-algebra and Z its center. Let \(\varphi : \mathrm{A} \otimes \mathrm{A}^{\circ} \to \operatorname{End}_{\mathrm{Z}}(\mathrm{A})\) be the homomorphism that sends \(a \otimes b\) to the endomorphism \(x \mapsto axb\) . Prove that \(\varphi\) is injective and that its image H is a quasi-simple algebra (use \(c\) ) and \(d\) ). Prove, moreover, that H is dense in \(\operatorname{End}_{\mathrm{Z}}(\mathrm{A})\) (VIII, p.~129, Exercise 3); we have \(\mathrm{H} = \operatorname{End}_{\mathrm{Z}}(\mathrm{A})\) if and only if A is finite-dimensional over Z (cf.~VIII, p.~128, Exercise 2).
\item
  Let A be a semisimple algebra with center K. Deduce from e) that the K-algebra \(A \otimes A^{o}\) is semisimple if and only if A has finite rank over K (cf.~VIII, p.~238, Theorem 2).
\end{enumerate}

\begin{enumerate}
\def\labelenumi{\arabic{enumi})}
\setcounter{enumi}{5}
\item
  Let A be a K-algebra with nonzero nilpotent radical. Prove that the homomorphism \(\varphi: A \otimes A^{\circ} \to \operatorname{End}_{K}(A)\) (Exercise 5) is not injective.
\item
  Let \(A_{1}\) and \(A_{2}\) be semisimple K-algebras and \(Z_{1}\) and \(Z_{2}\) their centers. Suppose that \(A_{1}\) has finite degree over K and that the ring \(Z_{1} \otimes Z_{2}\) is reduced. Let f be a K-algebra homomorphism from \(A_{2}\) to \(A_{1}\) . Prove that the commutant of \(f(\mathrm{A}_{2})\) in \(A_{1}\) is a semisimple K-algebra (view \(A_{1}\) as an \((\mathrm{A}_{2} \otimes \mathrm{A}_{1}^{\circ})\) -module and use Proposition 7 of VIII, p.~221).
\item
  Let K be a commutative field and L a radical extension of K of finite degree \textgreater{} 1. Let V be an infinite-dimensional vector space over L, and let A be the K-subalgebra of \(\operatorname{End}_{\mathrm{L}}(\mathrm{V})\) generated by the endomorphisms of finite rank and the identity mapping. Prove that the L-algebra \(\mathrm{A}_{(\mathrm{L})}\) has nonzero radical even though its center is equal to L (use Exercise 7).
\item
  Let A and B be K-algebras, M a free A-module, and N a B-module.
\end{enumerate}

\begin{enumerate}
\def\labelenumi{\alph{enumi})}
\item
  Prove that the image of the canonical homomorphism \(\vartheta\) from \(\mathrm{End}_{\mathrm{A}}(\mathrm{M})\otimes \mathrm{End}_{\mathrm{B}}(\mathrm{N})\) to \(\mathrm{End}_{\mathrm{A}\otimes \mathrm{B}}(\mathrm{M}\otimes \mathrm{N})\) is a dense subring of the ring \(\mathrm{End}_{\mathrm{A}\otimes \mathrm{B}}(\mathrm{M}\otimes \mathrm{N})\) .
\item
  If the ring of homotheties of the B-module N is dense in its bicommutant, then the same holds for the ring of homotheties of the \((A \otimes B)\) -module \(M \otimes N\) .
\item
  Suppose that there exist an element y of N and a sequence \((v_{n})\) of endomorphisms of N such that the elements \(v_{n}(y)\) generate an infinite-dimensional vector space over K. Prove that the mapping \(\vartheta\) is not surjective.
\end{enumerate}

\begin{enumerate}
\def\labelenumi{\arabic{enumi})}
\setcounter{enumi}{9}
\tightlist
\item
  Keep the assumptions of Exercise 9 and suppose, moreover, that the B-module N is simple; denote its commutant by D.
\end{enumerate}

\begin{enumerate}
\def\labelenumi{\alph{enumi})}
\item
  The mapping \(\vartheta\) is surjective if and only if M admits a finite basis over A or D is finite-dimensional over K (use Exercise 9, c)).
\item
  Prove that the \((\mathrm{A} \otimes \mathrm{B})\) -module \(M \otimes N\) is semisimple if and only if the ring \(A \otimes D^{o}\) is semisimple (reduce to the case \(M = A_{s}\) ; observe that \(M \otimes N\) is a free \(A \otimes D^{o}\) -module, and use Proposition 5 of VIII, p.~85 as well as Exercise 9, a)).
\end{enumerate}

¶ 11) Let A and B be K-algebras; suppose that the rings A and B are primitive and that their respective socles s and t (VIII, p.~130, Exercise 8) are not reduced to 0. Let D (resp. E) be the commutant of a minimal left ideal a of A (resp. b of B); suppose that the algebra D ⊗ E is simple. Prove that A ⊗ B is a primitive ring with socle s ⊗ t (observe that the (A ⊗ B)-module a ⊗ b is isotypical of finite length; use Exercise 9, a) of VIII, p.~131).

\begin{enumerate}
\def\labelenumi{\arabic{enumi})}
\setcounter{enumi}{11}
\tightlist
\item
  Let \(A_{1}\) be the K-algebra \(\mathrm{K}(\mathrm{X})[\mathrm{U}]_{1,d/d\mathrm{X}}\) (VIII, p.~11), \(A_{2}\) the K-algebra \(K[[T]]\) , and A the K-algebra \(A_{1} \otimes_{K} A_{2}\) . Identify \(A_{1}\) and \(A_{2}\) with subalgebras of A.
\end{enumerate}

\begin{enumerate}
\def\labelenumi{\alph{enumi})}
\item
  Prove that \(A_{2}\) is the center of A, that A does not contain any zero divisors, and that the invertible elements of A are those of \(A_{2}\) .
\item
  Prove that the two-sided ideals of A are the ideals \(AT^{n}\) for \(n \geqslant 0\) .
\item
  Prove that the ring A is primitive even though its center has nonzero radical (observe that a maximal left ideal of A containing 1 - TU does not contain any two-sided ideals other than 0).
\end{enumerate}

\begin{enumerate}
\def\labelenumi{\arabic{enumi})}
\setcounter{enumi}{12}
\tightlist
\item
  Let D and E be fields with centers containing K, and let V (resp. W) be a vector space over D (resp. E). Prove that the \((\mathrm{D} \otimes_{\mathrm{K}} \mathrm{E})\) -module \(V \otimes_{K} W\) is free.
\end{enumerate}

¶ 14) Let V be a vector space over a field D. Denote the endomorphism ring of the abelian group V by \(\Omega\) , and identify D with a subring of \(\Omega\) . Set \(A = \text{End}_{\text{D}}(\text{V})\) .

\begin{enumerate}
\def\labelenumi{\alph{enumi})}
\item
  Prove that A and D are linearly disjoint over their common center Z (cf.~Exercise 5, d)).
\item
  The image of \(D \otimes_{Z} A\) in \(\Omega\) is equal to \(\operatorname{End}_{\mathbb{Z}}(\mathrm{V})\) if and only if D has finite rank over Z. (To see that the condition is necessary, first deduce the case V = D from Exercise 5, e), and then fix an element \(x \neq 0\) of V and consider the subalgebra B of \(\operatorname{End}_{\mathbb{Z}}(\mathrm{V})\) consisting of the Z-endomorphisms u of V such that \(u(\mathrm{D}x) \subset \mathrm{D}x\) . Observe that the equality \(\operatorname{End}_{\mathbb{Z}}(\mathrm{V}) = \mathrm{DA}\) implies \(\mathrm{B} = \mathrm{D}(\mathrm{B} \cap \mathrm{A})\) .)
\end{enumerate}

\begin{enumerate}
\def\labelenumi{\arabic{enumi})}
\setcounter{enumi}{14}
\tightlist
\item
  \begin{enumerate}
  \def\labelenumii{\alph{enumii})}
  \tightlist
  \item
    Let \(\mathrm{K}\) be a commutative field, A a K-algebra, and \(\mathrm{B} = \mathrm{K}[\mathrm{X}_{\iota}]_{\iota \in \mathrm{I}}\) a polynomial algebra over \(\mathrm{K}\) . Prove that if every nonzero element of \(\mathrm{A}\) is regular, then the algebra \(\mathrm{A} \otimes_{\mathrm{K}} \mathrm{B}\) has the same property.
  \end{enumerate}
\end{enumerate}

\begin{enumerate}
\def\labelenumi{\alph{enumi})}
\setcounter{enumi}{1}
\tightlist
\item
  Deduce that if D is a field of finite rank over K and E is a purely transcendental extension of K, then the ring \(D \otimes_{K} E\) is a field.
\end{enumerate}

¶ 16) Let D be a field, V an infinite-dimensional right vector space over D, and A a dense subring of the ring \(\mathrm{End}_{\mathrm{D}}(\mathrm{V})\) whose socle \(\mathfrak{s}\) is not zero (VIII, p.~130, Exercises 8 and 9). Let K be a subfield of the center of D and B a K-algebra. Prove that if the radical of the K-algebra \(\mathrm{D}^{\circ} \otimes \mathrm{B}\) is nonzero, then so is that of \(\mathrm{A} \otimes \mathrm{B}\) . (Let \(\mathfrak{a}\) be the set of elements \(\sum u_{i} \otimes b_{i}\) of \(\mathfrak{s} \otimes \mathrm{B}\) satisfying \(\sum \langle u_{i}(x), x^{*} \rangle \otimes b_{i} \in \Re(\mathrm{D}^{\circ} \otimes \mathrm{B})\) for all \(x \in \mathrm{V}\) and \(x^{*} \in \mathrm{V}^{*}\) . Prove that \(\mathfrak{a}\) is a two-sided ideal of \(\mathrm{A} \otimes \mathrm{B}\) and, using Exercises 10, e) of VIII, p.~167 and 9 of VIII, p.~131, that it is nonzero and is contained in the radical of \(\mathrm{A} \otimes \mathrm{B}\) .)

\begin{enumerate}
\def\labelenumi{\arabic{enumi})}
\setcounter{enumi}{16}
\tightlist
\item
  Let K be a commutative field and D and E fields whose centers contain K. Suppose that \([E:K]=m\) is finite and that the center of one of the fields D and E is equal to K. The K-algebra \(D\otimes_{K}E\) is then isomorphic to the endomorphism algebra of a finite-dimensional vector space h over a field C (VIII, p.~222, Corollary 2).
\end{enumerate}

\begin{enumerate}
\def\labelenumi{\alph{enumi})}
\item
  Prove that \(h\) divides \(m\) (cf.~VIII, p.~205, formula (32)).
\item
  Let W be a vector space over D; then E is isomorphic to a K-subalgebra of \(\operatorname{End}_{\mathrm{D}}(\mathrm{W})\) if and only if the dimension of W is infinite or a multiple of m/h (note that W is then a \((\mathrm{D} \otimes_{\mathrm{K}} \mathrm{E})\) -module).
\item
  Prove that if \([\mathrm{D}:\mathrm{K}]\) is finite and prime to \(m\) , then \(\mathrm{D}\otimes_{\mathrm{K}}\mathrm{E}\) is a field (apply \(a\) ).
\end{enumerate}

\section{Absolutely Semisimple Algebras}

\subsection{Absolutely Semisimple Modules}

DEFINITION 1. --- Let K be a commutative field and A a K-algebra. We call an A-module M absolutely semisimple if the \(\mathrm{A}_{(\mathrm{L})}\) -module \(\mathrm{M}_{(\mathrm{L})}\) is semisimple for every extension L of K.

Every absolutely semisimple module is semisimple. Conversely, if the field K is perfect, then every semisimple A-module that is finite-dimensional over K is absolutely semisimple (VIII, p.~222, Corollary 1, a)) because every extension of a perfect field is separable (V, §7, No.~1, p.~36, Proposition 2).

PROPOSITION 1. --- Let K be a commutative field and A a K-algebra.

\begin{enumerate}
\def\labelenumi{\alph{enumi})}
\item
  Every direct sum of absolutely semisimple A-modules is an absolutely semisimple A-module. Every submodule and quotient module of an absolutely semisimple module is absolutely semisimple.
\item
  Let M be an A-module and let L be an extension of the field K. The A-module M is absolutely semisimple if and only if the \(\mathrm{A}_{(\mathrm{L})}\) -module \(\mathrm{M}_{(\mathrm{L})}\) is absolutely semisimple.
\end{enumerate}

Assertion a) follows from the analogous assertion for semisimple modules (VIII, p.~56, Corollaries 1 and 3).

Suppose that the A-module M is absolutely semisimple, and let \(L'\) be an extension of L. As an \(\mathrm{A}_{(\mathrm{L}')}\) -module, \(L' \otimes_{L} M_{(L)}\) is isomorphic to \(\mathrm{M}_{(\mathrm{L}')}\) (II, §5, No.~1, p.~273, Proposition 2); it is therefore a semisimple module. This proves that \(\mathrm{M}_{(\mathrm{L})}\) is absolutely semisimple.

Conversely, suppose that \(M_{(L)}\) is absolutely semisimple. Let \(L'\) be an extension of K. There exists a composite extension \((\Omega, u, v)\) of L and \(L'\) (V, §2, No.~4, p.~13, Corollary); we identify L and \(L'\) with subextensions of \(\Omega\) .

The \(A_{(\Omega)}\) -module \(M_{(\Omega)}\) is isomorphic to \((\mathrm{M}_{(\mathrm{L})})_{(\Omega)}\) ; it is therefore semisimple. But \(M_{(\Omega)}\) is also isomorphic to \((\mathrm{M}_{(\mathrm{L}^{\prime})})_{(\Omega)}\) , and Proposition 8, a) of VIII, p.~222 implies that \(M_{(L^{\prime})}\) is semisimple. So M is absolutely semisimple.

PROPOSITION 2. --- Let K be a commutative field, A a K-algebra, and M an A-module that is finite-dimensional over K. The following properties are equivalent:

\begin{enumerate}
\def\labelenumi{(\roman{enumi})}
\item
  The A-module M is absolutely semisimple.
\item
  There exists an extension \(\mathrm{P}\) of \(\mathrm{K}\) that is a perfect field such that the \(\mathrm{A}_{(\mathrm{P})}\) -module \(\mathrm{M}_{(\mathrm{P})}\) is semisimple.
\item
  The A-module M is semisimple, and the center of its commutant is an étale algebra over the field K.
\end{enumerate}

It is clear that (i) implies (ii). Assume that (ii) holds; then the \(A_{(P)}\) -module \(M_{(P)}\) is absolutely semisimple by Corollary 1, a) of VIII, p.~222. By Proposition 1, b), M is then absolutely semisimple. Therefore, (ii) implies (i).

Suppose that M is semisimple. Let Z be the center of the commutant of M; it is a commutative algebra of finite degree over the field K. If L is an extension of K, then it follows from Proposition 8, b) of VIII, p.~222 that the \(A_{(L)}\) -module \(M_{(L)}\) is semisimple if and only if the ring \(Z_{(L)}\) is reduced. The equivalence of (i) and (iii) therefore follows from Theorem 4 of V, §6, No.~7, p.~34.

\subsection{Algebras over Separably Closed Fields}

Lemma 1. --- Let D be a field and Z its center. Denote the characteristic exponent of Z by p.~Suppose that for every element a of D, there exists an integer \(m \geqslant 0\) such that \(a^{p^{m}}\) belongs to Z. Then the field D is commutative.

If p = 1, then D = Z. We therefore assume p \textgreater{} 1.

Let us give a proof by contradiction and suppose that D is not commutative. Let \(a\) be an element of \(\mathrm{D} - \mathrm{Z}\) and \(q\) a power of \(p\) such that \(a^q\) belongs to Z. Denote the identity mapping on D by I and the inner automorphism \(x \mapsto axa^{-1}\) from D to D associated with \(a\) by \(\sigma\) ; we have \(\sigma^q = \mathrm{I}\) because \(a^q\) belongs to Z. We have \(\sigma - \mathrm{I} \neq 0\) because \(a\) does not belong to Z, and we have \((\sigma - \mathrm{I})^q = \sigma^q - \mathrm{I} = 0\) because Z has characteristic \(p\) . Let \(f\) be the greatest natural number such that \((\sigma - \mathrm{I})^f \neq 0\) ; we have \(f \geqslant 1\) . Choose an element \(c\) of D such that \((\sigma - \mathrm{I})^f(c) \neq 0\) , and set

\[
x = (\sigma - \mathrm{I}) ^ {f - 1} (c), \qquad y = (\sigma - \mathrm{I}) (x) = (\sigma - \mathrm{I}) ^ {f} (c).
\]

By construction, we have \(y \neq 0\) and \(\sigma(y) = y\) ; if we set \(z = y^{-1}x\) , then it follows that

\[
\sigma (z) = \sigma (y) ^ {- 1} \sigma (x) = y ^ {- 1} (y + x) = 1 + z
\]

and therefore \(\sigma(z^{p^j}) = 1 + z^{p^j}\) for every natural number \(j\) . Choose an integer \(m \geqslant 0\) such that \(z^{p^m}\) belongs to the center Z of D; we have

\[
z ^ {p ^ {m}} = a z ^ {p ^ {m}} a ^ {- 1} = \sigma (z ^ {p ^ {m}}) = 1 + z ^ {p ^ {m}}.
\]

This contradiction implies Lemma 1.

PROPOSITION 3. --- Let K be a separably closed field (V, §7, No.~8, p.~45, Definition 4), and let D be an algebra of finite degree over K that is a field. Then D is commutative.

Denote the characteristic exponent of K by p.~Let a be an element of D. The ring \(K[a]\) is an algebraic extension of K (V, §3, No.~1, p.~17, Corollary 1). Since the field K is separably closed, it follows from V, §7, No.~7, p.~44, Proposition 13 that the algebra \(K[a]\) is a p-radical extension of K. Hence there exists an integer \(m \geqslant 0\) such that \(a^{p^{m}}\) belongs to K. By Lemma 1, the field D is consequently commutative.

COROLLARY. --- Let K be a separably closed field and A a semisimple algebra of finite degree over K. Then there exist an integer \(r \geqslant 0\) , strictly positive integers \(n_{1}, \ldots, n_{r}\) , and extensions \(K_{1}, \ldots, K_{r}\) of K of finite degree such that A is isomorphic to the algebra \(\prod_{i=1}^{r} \mathbf{M}_{n_{i}}(\mathbf{K}_{i})\) .

By the structure theorem for semisimple algebras (VIII, p.~135, Theorem 1), A is isomorphic to an algebra \(\prod_{i=1}^{r} \mathbf{M}_{n_i}(\mathrm{D}_i)\) , where \(r\) is an integer \(\geqslant 0\) , \(n_1, \ldots, n_r\) are strictly positive integers, and \(\mathrm{D}_1, \ldots, \mathrm{D}_r\) are K-algebras of finite degree that are fields. Since the field K is separably closed, each field \(\mathrm{D}_i\) is commutative by Proposition 3; the corollary follows.

\subsection{Absolutely Semisimple Algebras}

DEFINITION 2. --- Let K be a commutative field. We say that a K-algebra A is absolutely semisimple if the ring \(\mathrm{A}_{(\mathrm{L})}\) is semisimple for every extension L of K.

An absolutely semisimple algebra is semisimple. The K-algebra A is absolutely semisimple if and only if the A-module \(A_{s}\) is absolutely semisimple. By Proposition 1 of VIII, p.~229, we therefore obtain the following result: if

L is an extension of K, then the L-algebra \(A_{(L)}\) is absolutely semisimple if and only if the K-algebra A is absolutely semisimple.

THEOREM 1. --- Let K be a commutative field and A a K-algebra. The following properties are equivalent:

\begin{enumerate}
\def\labelenumi{(\roman{enumi})}
\item
  The K-algebra A is absolutely semisimple.
\item
  The algebra \(\mathrm{A}\) has finite degree over \(\mathrm{K}\) , and there exists an extension \(\mathrm{P}\) of \(\mathrm{K}\) that is a perfect field, for which the \(\mathrm{P}\) -algebra \(\mathrm{A}_{(\mathrm{P})}\) is semisimple.
\item
  The K-algebra A is semisimple and has finite degree over K, and its center is an étale K-algebra.
\item
  There exists a finite family \((n_{i},\mathrm{D}_{i})_{i\in\mathrm{I}}\) , where \(n_{i}\) is a strictly positive integer and \(D_{i}\) a K-algebra of finite degree that is a field, such that the center \(Z_{i}\) of \(D_{i}\) is a separable extension of K for every \(i\in I\) and A is isomorphic to the product of the matrix rings \(\mathbf{M}_{n_{i}}(\mathrm{D}_{i})\) .
\item
  There exist an extension \(\mathrm{L}\) of \(\mathrm{K}\) and a finite family of integers \((n_i)_{i\in \mathrm{I}}\) such that the L-algebra \(\mathrm{A}_{(\mathrm{L})}\) is isomorphic to the algebra \(\prod_{i\in \mathrm{I}}\mathbf{M}_{n_i}(\mathrm{L})\) .
\item
  There exist a Galois extension L of K of finite degree and a finite family of integers \((n_{i})\) such that \(\mathrm{A}_{(\mathrm{L})}\) is isomorphic to the product of the matrix algebras \(\mathbf{M}_{n_{i}}(\mathrm{L})\) .
\end{enumerate}

We first prove the implications (v) \(\Rightarrow\) (iv) \(\Rightarrow\) (iii) \(\Rightarrow\) (ii) \(\Rightarrow\) (i). Denote the center of A by Z.

If property (v) holds, then the L-algebra \(\mathrm{A}_{(\mathrm{L})}\) is semisimple and has finite degree over L, and its center (isomorphic to \(\mathrm{Z}_{(\mathrm{L})}\) , III, p.~41, Corollary) is isomorphic to \(L^{r}\) for some integer \(r \geqslant 0\) . By Corollary 2, a) of VIII, p.~222, the algebra A is semisimple and has finite degree over K. It is therefore isomorphic to a finite product of rings \(\prod_{i \in I} M_{n_{i}}(D_{i})\) with \(n_{i} \geqslant 1\) for every \(i \in I\) , where the field \(D_{i}\) is an algebra of finite degree over K. The center \(Z_{i}\) of \(D_{i}\) is a commutative field that is an extension of K, and Z is isomorphic to \(\prod_{i \in I} Z_{i}\) . Therefore, \(\mathrm{Z}_{(\mathrm{L})}\) is isomorphic, on the one hand, to \(\prod_{i \in I} (\mathrm{Z}_{i})_{(\mathrm{L})}\) and, on the other hand, to \(L^{r}\) . In other words, the algebra \(\prod_{i \in I} Z_{i}\) is étale over the field K (V, §6, No.~3, p.~28, Definition 1), and each of the extensions \(Z_{i}\) is separable over K (V, §6, No.~4, p.~31, Corollary, and V, §7, No.~1, p.~42). So (v) implies (iv).

If property (iv) holds, then it is clear that A is a semisimple algebra and has finite degree over K. Its center Z is isomorphic to the product \(\prod_{i\in I}Z_{i}\) of separable extensions of K of finite degree; it is therefore an étale algebra (loc. cit.). So (iv) implies (iii).

The implications (iii) \(\Rightarrow\) (ii) \(\Rightarrow\) (i) follow from Proposition 2 of VIII, p.~230 applied to the A-module \(\mathrm{A}_s\) .

Lemma 2. --- Let L be an algebraically closed field and D a field containing L in its center. If D is distinct from L, then there exists an extension \(L'\) of L such that the ring \(D \otimes_{L} L'\) is not left Artinian.

Let \(x\) be an element of \(\mathrm{D} - \mathrm{L}\) ; since \(\mathrm{L}\) is algebraically closed, the extension \(\mathrm{L}' = \mathrm{L}(x)\) of \(\mathrm{L}\) is not algebraic and \(x\) is transcendent over \(\mathrm{L}\) . The ring \(\mathrm{B} = \mathrm{L}' \otimes_{\mathrm{L}} \mathrm{L}'\) is then an integral domain by Proposition 5 of \(\mathrm{V}\) , §17, No.~4, p.~141.

The element \(y = x \otimes 1 - 1 \otimes x\) of B is not zero, but if \(\varphi\) is the homomorphism from B to \(L'\) that sends \(\xi \otimes \eta\) to \(\xi \eta\) , then we have \(\varphi(y) = 0\) , so y is not invertible in B. We view the ring \(C = D \otimes_{L} L'\) as a right module over its subring B; it is a free module because D is a right vector space over its subfield \(L'\) . Since y is a nonzero and noninvertible element of the integral domain B, right multiplication by y in C is a mapping \(R_y\) that is injective but not bijective. Now, \(R_y\) is an endomorphism of the left C-module \(C_s\) ; consequently (VIII, p.~28, Corollary 1), the ring C is not left Artinian.

Let us now prove that (i) implies (v). This is a consequence of the following lemma.

Lemma 3. --- Let A be an absolutely semisimple K-algebra and L an algebraically closed extension of K. Then the algebra \(\mathrm{A}_{(\mathrm{L})}\) is isomorphic to a product of finitely many matrix algebras over L.

The L-algebra \(\mathrm{A}_{(\mathrm{L})}\) is semisimple; it is therefore isomorphic to a product of finitely many algebras of the form \(\mathbf{M}_{n_{i}}(\mathrm{D}_{i})\) , where \(D_{i}\) is a field containing L in its center and \(n_{i}\) an integer \(\geqslant 1\) (VIII, p.~137, Remark 1).

Let \(L'\) be an extension of L. Since the K-algebra A is absolutely semisimple, the ring \(A_{(L')}\) is semisimple and therefore left Artinian. Now, the ring \(A_{(L')}\) is isomorphic to \(L' \otimes_L A_{(L)}\) , hence to \(\prod_{i \in I} M_{n_i}(L' \otimes_L D_i)\) ; by Proposition 5 of VIII, p.~7, each of the rings \(M_{n_i}(L' \otimes_L D_i)\) is therefore left Artinian.

Let \(n \geqslant 1\) be an integer and B a ring. Let \((\mathfrak{b}_r)_{r \geqslant 0}\) be a decreasing sequence of left ideals of B; denote by \(\mathfrak{c}_r\) the set of square matrices of order \(n\) with entries in \(\mathfrak{b}_r\) . Then \((\mathfrak{c}_r)_{r \geqslant 0}\) is a decreasing sequence of left ideals of \(\mathbf{M}_n(\mathrm{B})\) . In particular, if the ring \(\mathbf{M}_n(\mathrm{B})\) is left Artinian, then so is B.

By the above, for every \(i \in I\) and every extension \(L'\) of L, the ring \(D_{i} \otimes_{L} L'\) is left Artinian. By Lemma 2, we have \(D_{i} = L\) for every \(i \in I\) , which implies Lemma 3.

We will use the following lemma to prove the implication (i)⇒(vi).

Lemma 4. --- Let A and B be algebras over the field K that have finite generating sets, and let \(K'\) be an extension of K. If the \(K'\) -algebras \(\mathrm{A}_{(\mathrm{K}')}\) and \(\mathrm{B}_{(\mathrm{K}')}\) are isomorphic, then there exists a subextension L of \(K'\) , finitely generated over K, such that the L-algebras \(\mathrm{A}_{(\mathrm{L})}\) and \(\mathrm{B}_{(\mathrm{L})}\) are isomorphic.

Let \((e_{i})_{i\in\mathrm{I}}\) and \((f_{j})_{j\in\mathrm{J}}\) be finite generating sets for the algebras A and B, respectively. Let u be an isomorphism from \(\mathrm{A}_{(\mathrm{K}^{\prime})}\) to \(\mathrm{B}_{(\mathrm{K}^{\prime})}\) and v the inverse isomorphism; there exists a subextension L of \(K^{\prime}\) , finitely generated over K, such that we have \(u(1\otimes e_{i})\in\mathrm{B}_{(\mathrm{L})}\) for every \(i\in\mathrm{I}\) and \(v(1\otimes f_{j})\in\mathrm{A}_{(\mathrm{L})}\) for every \(j\in\mathrm{J}\) . Consequently, u sends \(\mathrm{A}_{(\mathrm{L})}\) into \(\mathrm{B}_{(\mathrm{L})}\) , and v sends \(\mathrm{B}_{(\mathrm{L})}\) into \(\mathrm{A}_{(\mathrm{L})}\) . The induced mappings \(u^{\prime}:\mathrm{A}_{(\mathrm{L})}\to\mathrm{B}_{(\mathrm{L})}\) and \(v^{\prime}:\mathrm{B}_{(\mathrm{L})}\to\mathrm{A}_{(\mathrm{L})}\) are ring homomorphisms; they are inverse bijections.

Let us complete the proof of the implication \((\mathrm{i})\Rightarrow(\mathrm{vi})\) . Let \(K'\) be a separable closure of K (V, §7, No.~8, p.~45). Then \(A_{(K')}\) is an absolutely semisimple algebra over \(K'\) . By the implication \((\mathrm{i})\Rightarrow(\mathrm{iv})\) , the \(K'\) -algebra \(A_{(K')}\) is isomorphic to a product \(\mathbf{M}_{n_{1}}(\mathrm{D}_{1})\times\cdots\times\mathbf{M}_{n_{r}}(\mathrm{D}_{r})\) , where \(D_{i}\) is a \(K'\) -algebra of finite degree that is a field whose center \(Z_{i}\) is a separable extension of \(K'\) . Since \(K'\) is separably closed, we have \(Z_{i}=K'\) . By Proposition 3 of VIII, p.~231, the field \(D_{i}\) is commutative. We therefore have \(D_{i}=K'\) . We denote the K-algebra \(\mathbf{M}_{n_{1}}(\mathrm{K})\times\cdots\times\mathbf{M}_{n_{r}}(\mathrm{K})\) by B. The \(K'\) -algebras \(A_{(K')}\) and \(B_{(K')}\) are isomorphic by the above. Every subextension of \(K'\) that is finitely generated over K is separable and has finite degree over K, hence is contained in a subextension L of \(K'\) that is Galois and of finite degree over K (V, §10, No.~1, p.~57, Proposition 2). The implication therefore follows from Lemma 4.

The implication (vi)⇒(v) is immediate.

COROLLARY 1. --- Let K be a commutative field, and let \(A_{1}\) and \(A_{2}\) be K-algebras. Suppose that \(A_{1}\) is absolutely semisimple.

\begin{enumerate}
\def\labelenumi{\alph{enumi})}
\item
  If \(\mathrm{A}_2\) is semisimple, then so is \(\mathrm{A}_1\otimes_{\mathrm{K}}\mathrm{A}_2\)
\item
  If \(\mathrm{A}_2\) is absolutely semisimple, then so is \(\mathrm{A}_1\otimes_{\mathrm{K}}\mathrm{A}_2\) .
\end{enumerate}

Denote the center of \(A_{1}\) by \(Z_{1}\) and that of \(A_{2}\) by \(Z_{2}\) . The center Z of \(A_{1} \otimes_{K} A_{2}\) is equal to \(Z_{1} \otimes_{K} Z_{2}\) by the corollary of III, §4, No.~4, p.~468. Suppose that \(A_{2}\) is semisimple; then \(Z_{2}\) is a reduced algebra (VIII, p.~137, Propositions 2 and 3). By Theorem 1, \(Z_{1}\) is an étale and therefore separable

K-algebra. By Proposition 5 of V, §15, No.~2, p.~120, the ring \(Z = Z_{1} \otimes_{K} Z_{2}\) is reduced; since \(A_{1}\) has finite degree over K (Theorem 1), it follows from Proposition 7 of VIII, p.~221 that the ring \(A_{1} \otimes_{K} A_{2}\) is semisimple.

Now suppose that \(A_{2}\) is absolutely semisimple. Let L be an extension of K. Then the algebra \(\mathrm{A}_{1(\mathrm{L})}\) is absolutely semisimple, and the algebra \(\mathrm{A}_{2(\mathrm{L})}\) is semisimple. Therefore, by a), the algebra \(\mathrm{A}_{1} \otimes_{\mathrm{K}} \mathrm{A}_{2(\mathrm{L})}\) is semisimple.

COROLLARY 2. --- Let K be a separably closed field, and let A be an absolutely semisimple K-algebra. Then there exist an integer \(r \geqslant 0\) and strictly positive integers \(n_{1}, \ldots, n_{r}\) such that the algebra A is isomorphic to the algebra \(\prod_{i=1}^{r} \mathbf{M}_{n_{i}}(\mathrm{K})\) .

By Theorem 1, A is isomorphic to an algebra of the form \(\prod_{i=1}^{r} \mathbf{M}_{n_i}(\mathrm{D}_i)\) for an integer \(r \geqslant 0\) , integers \(n_1, \ldots, n_r\) , and K-algebras of finite degree \(\mathrm{D}_1, \ldots, \mathrm{D}_r\) that are fields and whose centers are separable extensions of K and therefore equal to K. By Proposition 3 of VIII, p.~231, we have \(\mathrm{D}_i = \mathrm{K}\) for \(i \in [1, r]\) .

Example. --- A commutative K-algebra is absolutely semisimple if and only if it is étale: this follows from the definition (V, §6, No.~3, p.~28, Definition 1) and the equivalence of properties (i) and (v) of Theorem 1.

\subsection{Characterization of Absolutely Semisimple Modules}

PROPOSITION 4. --- Let K be a commutative field and A a K-algebra.

\begin{enumerate}
\def\labelenumi{\alph{enumi})}
\item
  Let M be a semisimple A-module. The A-module M is absolutely semisimple if and only if every simple module belonging to the support of M is.
\item
  Let S be a simple A-module, and let D be its commutant. The following properties are equivalent:
\end{enumerate}

\begin{enumerate}
\def\labelenumi{(\roman{enumi})}
\item
  The A-module S is absolutely semisimple.
\item
  The K-algebra \(\mathbf{D}\) is absolutely semisimple.
\item
  The K-algebra D is a field, has finite degree over K, and its center is a separable extension of K.
\end{enumerate}

Assertion a) follows from Proposition 1, a) of VIII, p.~229. Let S and D be as in b), and let L be an extension of K. The \(A_{(L)}\) -module \(S_{(L)}\) is semisimple if and only if the ring \(D_{(L)}\) is (VIII, p.~222, Proposition 8, c)). This proves the equivalence of (i) and (ii), and that of (ii) and (iii) follows from Theorem 1 because D is a field.

COROLLARY. --- Let K be a commutative field, let \(A_{1}\) and \(A_{2}\) be K-algebras, and let \(M_{1}\) be an absolutely semisimple \(A_{1}\) -module and \(M_{2}\) a semisimple \(A_{2}\) -module. Then \(M_{1} \otimes_{K} M_{2}\) is a semisimple module over the ring \(A_{1} \otimes_{K} A_{2}\) .

The module \(M_{1}\) is the direct sum of absolutely semisimple simple \(A_{1}\) -modules (Proposition 4). It therefore suffices to prove the assertion in the case when the modules \(M_{1}\) and \(M_{2}\) are simple. Denote their commutants by \(D_{1}\) and \(D_{2}\) . The K-algebra \(D_{1}\) is absolutely semisimple (loc. cit.); by Corollary 1 of VIII, p.~234, the K-algebra \(D_{1} \otimes_{K} D_{2}\) is semisimple. It then follows from Corollary 1 of VIII, p.~215 that the \((\mathrm{A}_{1} \otimes_{\mathrm{K}} \mathrm{A}_{2})\) -module \(M_{1} \otimes_{K} M_{2}\) is semisimple.

\subsection{Derivations on Semisimple Algebras}

In this subsection and the next ones, K is a commutative ring, A a K-algebra, B the K-algebra \(A \otimes_{K} A^{o}\) , and \(\varepsilon\) the K-linear mapping from B to A defined by \(\varepsilon(x \otimes y) = xy\) for x, y in A.

Recall (III, §4, No.~3, p.~467) that every (A, A)-bimodule P can be viewed as a left B-module, with external law given by \((a \otimes a')p = apa'\) for \(a, a'\) in A and p in P. Conversely, every B-module can be viewed as an (A, A)-bimodule. We endow A with its canonical (A, A)-bimodule structure and with the corresponding B-module structure; we endow B with the (A, A)-bimodule structure corresponding to the B-module \(B_s\) . We therefore have

\[
a (x \otimes y) a ^ {\prime} = (a \otimes a ^ {\prime}) (x \otimes y) = a x \otimes y a ^ {\prime}
\]

for \(a, a', x, y\) in A, where the product \(ya'\) is calculated in the algebra A.

The K-linear mapping \(\varepsilon\) is a homomorphism of (A, A)-bimodules.

PROPOSITION 5. --- The following properties are equivalent:

\begin{enumerate}
\def\labelenumi{(\roman{enumi})}
\item
  The B-module A is projective.
\item
  There exists an element \(e\) of the (A,A)-bimodule B satisfying the following two conditions: \(\varepsilon(e) = 1\) and \(ae = ea\) for every \(a \in \mathrm{A}\) .
\end{enumerate}

The mapping \(\varepsilon: \mathrm{B} \to \mathrm{A}\) is surjective because we have \(\varepsilon(a \otimes 1) = a\) for every \(a \in \mathrm{A}\) ; it is a homomorphism of (A, A)-modules, hence a B-linear mapping. If the B-module A is projective, then there exists a section \(s\) of \(\varepsilon\) (II, §2, No.~2, p.~231, Proposition 4); it is a homomorphism of (A, A)-bimodules from A to B. If we set \(e = s(1)\) , then we have \(\varepsilon(e) = \varepsilon(s(1)) = 1\) and \(ae = s(a) = ea\) for every \(a \in \mathrm{A}\) . So (i) implies (ii).

Conversely, let e be an element of B satisfying the conditions in item (ii). Define a mapping s from A to B by the formula

\[
s (a) = a e = e a.\tag{1}
\]

It is a homomorphism of (A, A)-modules, and we have \(\varepsilon \circ s = 1_{\mathrm{A}}\) ; in other words, \(s\) is a B-linear section of the surjective mapping \(\varepsilon\) . Consequently, the B-module A is isomorphic to the direct factor submodule \(s(\mathrm{A})\) of \(\mathrm{B}_s\) (II, §1, No.~9, p.~211, Proposition 15) and is therefore projective (II, §2, No.~2, p.~231, Proposition 4). This proves that (ii) implies (i).

Remarks. --- 1) Let \(e = \sum_{i=1}^{r} a_i \otimes a_i'\) be an element of B. The conditions in item (ii) of Proposition 5 translate to the formulas

\[
\sum_ {i = 1} ^ {r} a _ {i} a _ {i} ^ {\prime} = 1,\tag{2}
\]

\[
\sum_ {i = 1} ^ {r} a a _ {i} \otimes a _ {i} ^ {\prime} = \sum_ {i = 1} ^ {r} a _ {i} \otimes a _ {i} ^ {\prime} a \quad \text { for   every } a \in A.\tag{3}
\]

When they are satisfied, e is an idempotent in B. Indeed, we then have the relations

\[
e ^ {2} = \sum_ {i = 1} ^ {r} a _ {i} e a _ {i} ^ {\prime} = \sum_ {i = 1} ^ {r} e a _ {i} a _ {i} ^ {\prime} = e.
\]

\begin{enumerate}
\def\labelenumi{\arabic{enumi})}
\setcounter{enumi}{1}
\tightlist
\item
  Let K be a commutative field, let A be a K-algebra, and let M be an A-module. The group \(\operatorname{End}_{\mathrm{K}}(\mathrm{M})\) is endowed with an (A, A)-bimodule structure defined by
\end{enumerate}

\[
a u a ^ {\prime} (x) = a u \left(a ^ {\prime} x\right)
\]

for every \(a, a' \in A\) , every \(u \in \operatorname{End}_{\mathrm{K}}(\mathrm{M})\) , and every \(x \in M\) . We endow it with the associated B-module structure. Let \(e = \sum_{i=1}^{r} a_i \otimes a'_i\) be an element of B satisfying the conditions in part (ii) of Proposition 5, so also relations (2) and (3). If \(p \in \operatorname{End}_{\mathrm{K}}(\mathrm{M})\) is a projector whose image N is an A-submodule of M, then ep is an A-linear projector with the same image.

Indeed, the image of \(ep\) is contained in N. If \(x\) belongs to N, then so does \(a_i' x\) , so that we have \(p(a_i' x) = a_i' x\) and

\[
e p (x) = \sum_ {i = 1} ^ {r} a _ {i} a _ {i} ^ {\prime} x = x
\]

by formula (2). We deduce from formula (3) that \(aep(x) = ep(ax)\) for every \(a \in \mathrm{A}\) and every \(x \in \mathrm{M}\) , which proves that \(ep\) is A-linear.

THEOREM 2. --- Let K be a commutative field and A a K-algebra. The following properties are equivalent:

\begin{enumerate}
\def\labelenumi{(\roman{enumi})}
\item
  The K-algebra A is absolutely semisimple.
\item
  The K-algebra \(\mathrm{B} = \mathrm{A}\otimes_{\mathrm{K}}\mathrm{A}^{\circ}\) is semisimple.
\item
  The B-module A is projective.
\item
  There exists an element \(e\) of the (A,A)-bimodule B satisfying \(\varepsilon(e) = 1\) and \(ae = ea\) for every \(a \in \mathrm{A}\) .
\end{enumerate}

Suppose that the algebra A is absolutely semisimple, hence semisimple. Then the algebra \(A^{o}\) is semisimple (VIII, p.~137, Proposition 2), and it follows from Corollary 1 of VIII, p.~234 that \(B = A \otimes_{K} A^{o}\) is a semisimple K-algebra. Therefore, (i) implies (ii).

Since every module over a semisimple ring is projective (VIII, p.~138, Proposition 4), (ii) implies (iii).

The equivalence of (iii) and (iv) follows from Proposition 5. To complete the proof, let us show that (iv) implies (i). Let \(e = \sum_{i=1}^{r} a_i \otimes a_i'\) be an element of B satisfying the conditions in item (ii) of Proposition 5. Let L be an extension of the field K; we must prove that the ring \(\mathrm{A}_{(\mathrm{L})}\) is semisimple or, equivalently, that every \(\mathrm{A}_{(\mathrm{L})}\) -module is semisimple (VIII, p.~138, Proposition 4). Let M be an \(\mathrm{A}_{(\mathrm{L})}\) -module, and let N be a submodule of M; we view M as a left (A,L)-bimodule and N as a sub-bimodule (III, §4, No.~3, p.~466). Since L is a field, there exists an L-linear projector \(u\) in M with image N. Since the homotheties \(a_{\mathrm{M}}\) associated with the elements \(a\) of A are L-linear, there exists a unique group homomorphism from \(\mathrm{A} \otimes_{\mathrm{K}} \mathrm{A}^{\circ}\) to \(\operatorname{End}_{\mathrm{L}}(\mathrm{M})\) that sends an element \(a \otimes a'\) to the L-linear mapping \(x \mapsto au(a'x)\) . We denote the image of \(e\) by this homomorphism by \(v\) ; it follows from Remark 2 that \(v\) is an \(\mathrm{A}_{(\mathrm{L})}\) -linear projector with image N. The kernel of \(v\) is an \(\mathrm{A}_{(\mathrm{L})}\) -submodule of M, supplementary to N. By Corollary 2 of VIII, p.~56, the \(\mathrm{A}_{(\mathrm{L})}\) -module M is semisimple.

Remark 3. --- We know (VIII, p.~234, Corollary 1) that the tensor product of two absolutely semisimple algebras over a commutative field is absolutely semisimple. Consequently, if the algebra A is absolutely semisimple, then so is the algebra \(B = A \otimes_{K} A^{\circ}\) .

\subsection{Cohomology of Algebras}

In this subsection, K is a commutative ring, A a K-algebra, B the K-algebra \(A \otimes_{K} A^{\circ}\) , and \(\varepsilon\) the K-linear mapping from B to A defined by \(\varepsilon(x \otimes y) = xy\) for x, y in A. For \(n \in N\) , we denote the tensor product over K of \(n + 2\) copies of the K-module A by \(B_{n}\) . We view it as an (A, A)-bimodule (and also as a B-module); we endow it with the structure of a left A-module deduced from the left A-module structure of the first factor of the tensor product and with the structure of a right A-module deduced from the right A-module structure of the last factor. In particular, \(B_{0}\) is just the (A, A)-bimodule B.

For every integer \(n \geqslant 1\) , we define homomorphisms of bimodules \(d_{n}^{i}\) for \(i \in [0, n]\) and \(d_{n}\) from \(B_{n}\) to \(B_{n-1}\) by the formulas

\[
d _ {n} ^ {i} (x _ {0} \otimes \dots \otimes x _ {n + 1}) = x _ {0} \otimes \dots \otimes x _ {i} x _ {i + 1} \otimes \dots \otimes x _ {n + 1}\tag{4}
\]

for \(i\in [0,n]\) and

\[
d _ {n} = \sum_ {i = 0} ^ {n} (- 1) ^ {i} d _ {n} ^ {i}.\tag{5}
\]

We denote the mapping \(\varepsilon : \mathrm{B}_0 \to \mathrm{A}\) by \(d_0 = d_0^0\) .

Let \(n\) be an integer \(\geqslant 1\) . For \(0 \leqslant i < j \leqslant n\) , we have

\[
d _ {n - 1} ^ {i} \circ d _ {n} ^ {j} = d _ {n - 1} ^ {j - 1} \circ d _ {n} ^ {i},\tag{6}
\]

and from this we deduce

\[
\begin{array}{c} d _ {n - 1} \circ d _ {n} = \sum_ {0 \leqslant i <   j \leqslant n} (- 1) ^ {i + j} d _ {n - 1} ^ {i} \circ d _ {n} ^ {j} + \sum_ {0 \leqslant j \leqslant i \leqslant n - 1} (- 1) ^ {i + j} d _ {n - 1} ^ {i} \circ d _ {n} ^ {j} \\ = \sum_ {0 \leqslant i <   j \leqslant n} (- 1) ^ {i + j} d _ {n - 1} ^ {j - 1} \circ d _ {n} ^ {i} + \sum_ {0 \leqslant j \leqslant i \leqslant n - 1} (- 1) ^ {i + j} d _ {n - 1} ^ {i} \circ d _ {n} ^ {j} \end{array}
\]

and consequently

\[
d _ {n - 1} \circ d _ {n} = 0.\tag{7}
\]

Let P be an (A,A)-bimodule. For any integer \(n \geqslant 0\) , we denote the K-module of K-multilinear mappings from \(A^{n}\) to P by \(C^{n}(A,P)\) . The mapping \(\alpha^{n}: C^{n}(A,P) \to \operatorname{Hom}_{\mathrm{B}}(\mathrm{B}_{n},\mathrm{P})\) that sends \(f \in \mathrm{C}^{n}(A,P)\) to the homomorphism \(\alpha^{n}(f)\) characterized by

\[
\alpha^ {n} (f) \left(x _ {0} \otimes \dots \otimes x _ {n + 1}\right) = x _ {0} f (x _ {1}, \ldots , x _ {n}) x _ {n + 1}\tag{8}
\]

is an isomorphism of K-modules.

We denote by \(\partial^{n}\) (for \(n \geqslant 0\) ) the unique K-linear mapping from \(\mathrm{C}^{n}(\mathrm{A},\mathrm{P})\) to \(\mathrm{C}^{n+1}(\mathrm{A},\mathrm{P})\) that makes the following diagram commutative:

\[
\begin{array}{c} \mathrm{C} ^ {n} (\mathrm{A}, \mathrm{P}) \xrightarrow {\partial^ {n}} \mathrm{C} ^ {n + 1} (\mathrm{A}, \mathrm{P}) \\ \Biggl \downarrow_ {\alpha^ {n}} \\ \mathrm{Hom} _ {\mathrm{B}} (\mathrm{B} _ {n}, \mathrm{P}) \xrightarrow {\mathrm{Hom} (d _ {n + 1} , 1 _ {\mathrm{P}})} \mathrm{Hom} _ {\mathrm{B}} (\mathrm{B} _ {n + 1}, \mathrm{P}). \end{array}
\]

By definition, we therefore have

\[
(\alpha^ {n + 1} \circ \partial^ {n}) (f) = \alpha^ {n} (f) \circ d _ {n + 1}\tag{9}
\]

for every \(f \in \mathrm{C}^n(\mathrm{A},\mathrm{P})\) . In other words, we have

\[
\partial^ {n} (f) \left(x _ {0}, \dots , x _ {n}\right) = \alpha^ {n} (f) \left(d _ {n + 1} \left(1 \otimes x _ {0} \otimes \dots \otimes x _ {n} \otimes 1\right)\right)
\]

for \(x_0, \ldots, x_n\) in A and \(f\) in \(\mathrm{C}^n(\mathrm{A}, \mathrm{P})\) , that is,

\[
\begin{array}{r l} & {\partial^ {n} (f) (x _ {0}, \ldots , x _ {n}) = x _ {0} f (x _ {1}, \ldots , x _ {n})} \\ & {\qquad + \sum_ {i = 0} ^ {n - 1} (- 1) ^ {i + 1} f (x _ {0}, \ldots , x _ {i - 1}, x _ {i} x _ {i + 1}, x _ {i + 2}, \ldots , x _ {n})} \\ & {\qquad \qquad \qquad \qquad \qquad + (- 1) ^ {n + 1} f (x _ {0}, \ldots , x _ {n - 1}) x _ {n}.} \end{array}\tag{10}
\]

By (7) and (9), we have

\[
\partial^ {n + 1} \circ \partial^ {n} = 0 \quad \text {   for   every   } n \geqslant 0.\tag{11}
\]

We denote the K-module \(\operatorname{Ker} \partial^0\) by \(\mathrm{H}^0(\mathrm{A},\mathrm{P})\) and, for \(n \geqslant 1\) , the K-module \(\operatorname{Ker} \partial^n/\operatorname{Im} \partial^{n-1}\) by \(\mathrm{H}^n(\mathrm{A},\mathrm{P})\) . We identify the K-module \(\mathrm{C}^0(\mathrm{A},\mathrm{P})\) with \(\mathrm{P}\) , and we have \(\mathrm{C}^1(\mathrm{A},\mathrm{P}) = \operatorname{Hom}_\mathrm{K}(\mathrm{A},\mathrm{P})\) . The mappings \(\partial^n\) for \(n \leqslant 2\) are given by the formulas

\[
\partial^ {0} (p) (a) = a p - p a \qquad \mathrm{for\ every} p \in \mathrm{P},\tag{12}
\]

\[
\partial^ {1} (f) \left(a, a ^ {\prime}\right) = a f \left(a ^ {\prime}\right) - f \left(a a ^ {\prime}\right) + f (a) a ^ {\prime} \quad \text { for } f \in \mathrm{C} ^ {1} (\mathrm{A}, \mathrm{P}),\tag{13}
\]

\[
\partial^ {2} (f) \left(a, a ^ {\prime}, a ^ {\prime \prime}\right) = a f \left(a ^ {\prime}, a ^ {\prime \prime}\right) - f \left(a a ^ {\prime}, a ^ {\prime \prime}\right) + f \left(a, a ^ {\prime} a ^ {\prime \prime}\right) - f \left(a, a ^ {\prime}\right) a ^ {\prime \prime}\tag{14}
\]

for \(f \in \mathrm{C}^2(\mathrm{A},\mathrm{P})\) .

So \(H^{0}(A,P)\) is the K-submodule of P consisting of the elements p such that ap = pa for every \(a \in A\) , and \(H^{1}(A,P)\) is the quotient of the K-module \(\mathrm{Der}_{\mathrm{K}}(\mathrm{A},\mathrm{P})\) of K-derivations from A to P (III, §10, No.~2, p.~553) by the K-submodule consisting of the derivations of the form \(a \mapsto ap - pa\) with \(p \in P\) (called inner derivations).

\subsection{Cohomology of Absolutely Semisimple Algebras}

PROPOSITION 6. --- Let K be a commutative ring and A a K-algebra. Let \(e = \sum_{i=1}^{r} a_i \otimes a_i'\) be an element of \(B = A \otimes_K A^o\) satisfying the conditions in item (ii) of Proposition 5 of VIII, p.~236. For any integer \(n \geqslant 1\) and any element f of \(C^n(A, P)\) , we denote by \(\gamma^n(f)\) the element of \(C^{n-1}(A, P)\) defined by the formula

\[
\gamma^ {n} (f) \left(x _ {1}, \dots , x _ {n - 1}\right) = \sum_ {i = 1} ^ {r} a _ {i} f \left(a _ {i} ^ {\prime}, x _ {1} \dots , x _ {n - 1}\right).\tag{15}
\]

We then have

\[
\partial^ {n - 1} (\gamma^ {n} (f)) + \gamma^ {n + 1} (\partial^ {n} (f)) = f\tag{16}
\]

for every integer \(n \geqslant 1\) and every \(f \in \mathrm{C}^n(\mathrm{A},\mathrm{P})\) .

*Remark 1. --- The morphisms \(\partial^n: \mathrm{C}^n(\mathrm{A},\mathrm{P}) \to \mathrm{C}^{n+1}(\mathrm{A},\mathrm{P})\) define a complex \((\mathrm{C}(\mathrm{A},\mathrm{P}),\partial)\) of K-modules (X, §2, n° 1, p.~24). The mapping \(\gamma_n\) therefore defines a homotopy from this complex to itself linking 0 to \(\mathrm{Id}_{\mathrm{C}(\mathrm{A},\mathrm{P})}\) (X, §2, n° 4, p.~32, définition 4). *

We keep the notation of No.6. For every integer \(n \geqslant 0\) , we define a mapping \(h_n: \mathrm{B}_n \to \mathrm{B}_{n+1}\) by the formula

\[
h _ {n} (x) = d _ {n + 2} ^ {1} (e \otimes x) = \sum_ {i = 1} ^ {r} a _ {i} \otimes a _ {i} ^ {\prime} x.
\]

It is a homomorphism of \((A, A)\) -bimodules (formula (3)).

Lemma 5. --- We have the relation

\[
d _ {n + 1} \circ h _ {n} + h _ {n - 1} \circ d _ {n} = 1 _ {\mathrm{B} _ {n}}\tag{17}
\]

for every \(n \geqslant 1\) .

Let \(x\in \mathrm{B}_n\) ; we have

\[
\begin{array}{c} (d _ {n + 1} \circ h _ {n}) (x) = (d _ {n + 1} \circ d _ {n + 2} ^ {1}) (e \otimes x) \\ = (d _ {n + 1} ^ {0} \circ d _ {n + 2} ^ {1}) (e \otimes x) - \sum_ {i = 2} ^ {n + 2} (- 1) ^ {i} (d _ {n + 1} ^ {i - 1} \circ d _ {n + 2} ^ {1}) (e \otimes x), \end{array}
\]

so, by formula (6),

\[
(d _ {n + 1} \circ h _ {n}) (x) = (d _ {n + 1} ^ {0} \circ d _ {n + 2} ^ {0}) (e \otimes x) - \sum_ {i = 2} ^ {n + 2} (- 1) ^ {i} (d _ {n + 1} ^ {1} \circ d _ {n + 2} ^ {i}) (e \otimes x).
\]

But we have

\[
(d _ {n + 1} ^ {0} \circ d _ {n + 2} ^ {0}) (e \otimes x) = \varepsilon (e) x = x
\]

by property (ii) of Proposition 5 of VIII, p.~236 and, for \(i \geqslant 2\) ,

\[
d _ {n + 2} ^ {i} (e \otimes x) = e \otimes d _ {n} ^ {i - 2} (x),
\]

which gives

\[
(d _ {n + 1} \circ h _ {n}) (x) = x - d _ {n + 1} ^ {1} (e \otimes d _ {n} (x)) = x - h _ {n - 1} \circ d _ {n} (x)
\]

and therefore formula (17).

We can finish the proof of Proposition 6 using Lemma 5. Let n be an integer \(\geqslant 1\) and f an element of \(C^{n}(A,P)\) . By construction, we have

\[
\alpha^ {n - 1} (\gamma^ {n} (f)) = \alpha^ {n} (f) \circ h _ {n - 1}\tag{18}
\]

and, consequently, by formulas (9) and (18),

\[
\begin{array}{r l} \alpha^ {n} (\partial^ {n - 1} (\gamma^ {n} (f)) + \gamma^ {n + 1} (\partial^ {n} (f))) & = \alpha^ {n - 1} (\gamma^ {n} (f)) \circ d _ {n} + \alpha^ {n + 1} (\partial^ {n} (f)) \circ h _ {n} \\ & = \alpha^ {n} (f) \circ h _ {n - 1} \circ d _ {n} + \alpha^ {n} (f) \circ d _ {n + 1} \circ h _ {n} \\ & = \alpha^ {n} (f), \end{array}
\]

where the last equality follows from (17). Since \(\alpha^{n}\) is bijective, the proposition follows.

THEOREM 3. --- Let K be a commutative ring, A a K-algebra, and P an (A,A)-bimodule. Suppose that the \((\mathrm{A}\otimes_{\mathrm{K}}\mathrm{A}^{\circ})\) -module A is projective. We then have \(\mathrm{H}^{n}(\mathrm{A},\mathrm{P})=0\) for every integer \(n\geqslant1\) .

We must prove that for every integer \(n \geqslant 1\) , every element f of \(C^{n}(A, P)\) with \(\partial^{n}(f) = 0\) is of the form \(\partial^{n-1}(g)\) for an element g of \(C^{n-1}(A, P)\) . By Proposition 5, this is an immediate consequence of Proposition 6.

COROLLARY. --- Every K-derivation from A to P is inner.

This is a translation of the equality \(H^{1}(A,P)=0\) .

Remarks. --- 2) The assumptions of Theorem 3 are, in particular, satisfied when K is a field and A an absolutely semisimple K-algebra (VIII, p.~238, Theorem 2).

*3) Suppose that the K-module A is projective. Theorem 3 can also be proved as follows. The complex \((\bigoplus_{n\geqslant 0}\mathrm{B}_n,d)\) and the homomorphism \(\varepsilon :\mathrm{B}_0\to \mathrm{A}\) define a projective resolution of the B-module A; therefore, for every \(n\geqslant 0\) , the K-module \(\mathrm{H}^n (\mathrm{A},\mathrm{P})\) is isomorphic to \(\operatorname {Ext}_{\mathrm{B}}^{n}(\mathrm{A},\mathrm{P})\) (X, §6, n° 1, p.~100, théorème 1). If the B-module A is projective, then the K-modules \(\operatorname {Ext}_{\mathrm{B}}^{n}(\mathrm{A},\mathrm{P})\) are zero for \(n\geqslant 1\) (X, §5, n° 3, p.~88, corollaire de la proposition 5), which implies that \(\mathrm{H}^n (\mathrm{A},\mathrm{P})\) is zero. Conversely, if \(\mathrm{H}^1 (\mathrm{A},\mathrm{P})\) is zero for every (A, A)-bimodule P, then the B-module A is projective (X, §5, n° 5, p.~93, proposition 10).*

\subsection{The Splitting of Artinian Algebras}

In this subsection, K is a commutative ring and A a K-algebra. Let r be the radical of A. We denote the quotient algebra \(A/r\) by \(\overline{A}\) and the canonical mapping from A to \(\overline{A}\) by \(\pi\) . We are interested in the subalgebras S of A such that \(A = S \oplus r\) .

We denote by \(\Sigma\) the set of K-linear sections \(s\) of \(\pi\) satisfying \(s(\alpha\beta) = s(\alpha)s(\beta)\) for \(\alpha, \beta\) in \(\overline{\mathrm{A}}\) . Note that such a section necessarily satisfies \(s(1) = 1\) (in other words, \(s\) is a ring homomorphism): we indeed have \(s(1)^2 = s(1)\) , and \(s(1)\) is invertible because it belongs to \(1 + \mathfrak{r}\) (VIII, p.~156, Theorem 1). If \(s\) is an element of \(\Sigma\) , then the image S of \(s\) is a subalgebra of A, and we have \(\mathrm{A} = \mathrm{S} \oplus \mathfrak{r}\) . Conversely, if S is a subalgebra of A such that \(\mathrm{A} = \mathrm{S} \oplus \mathfrak{r}\) , then the restriction of \(\pi\) to S is bijective, and the inverse bijection defines an element of \(\Sigma\) with image S.

By Jacobson's theorem (loc. cit.), every element of \(1 + \mathfrak{r}\) is invertible in A. We call an inner automorphism of A of the form \(a \mapsto xax^{-1}\) with \(x \in 1 + \mathfrak{r}\) a special automorphism.

PROPOSITION 7. --- Suppose that the \((\overline{\mathrm{A}}\otimes_{\mathrm{K}}\overline{\mathrm{A}}^{\mathrm{o}})\) -module \(\overline{\mathrm{A}}\) is projective.

\begin{enumerate}
\def\labelenumi{\alph{enumi})}
\item
  Let \(S_1\) and \(S_2\) be subalgebras of A satisfying \(A = S_1 \oplus \mathfrak{r} = S_2 \oplus \mathfrak{r}\) . There exists a special automorphism of A transforming \(S_1\) into \(S_2\) .
\item
  Suppose that \(\pi\) has a K-linear section and that the radical r of A is nilpotent. Then there exists a subalgebra S of A satisfying \(A = S \oplus r\) .
\end{enumerate}

Let \(S_{1}\) and \(S_{2}\) be as in a). Let \(s_{1}\) and \(s_{2}\) be the elements of the set \(\Sigma\) corresponding to the subalgebras \(S_{1}\) and \(S_{2}\) . Let \(\varepsilon\) be the K-linear mapping from \(A \otimes_{K} A\) to \(A\) given by \(\varepsilon(a \otimes b) = ab\) . By Proposition 5 of VIII, p.~236 and Remark 1 of VIII, p.~237, there exists an element \(e = \sum_{i=1}^{r} \alpha_{i} \otimes \alpha_{i}'\) of \(\overline{A} \otimes_{K} \overline{A}\) satisfying \(\sum_{i=1}^{r} \alpha_{i} \alpha_{i}' = 1\) and \(\sum_{i=1}^{r} \alpha \alpha_{i} \otimes \alpha_{i}' = \sum_{i=1}^{r} \alpha_{i} \otimes \alpha_{i}' \alpha\) for every \(\alpha \in \overline{A}\) . Set \(x = \sum_{i=1}^{r} s_{1}(\alpha_{i}) s_{2}(\alpha_{i}')\) . We have \(\pi(x) = \sum_{i=1}^{r} \alpha_{i} \alpha_{i}' = 1\) and therefore \(x \in 1 + r\) . Let \(\alpha\) be an element of \(\overline{A}\) . We have

\[
\begin{array}{l} s _ {1} (\alpha) x = \sum_ {i = 1} ^ {r} s _ {1} (\alpha \alpha_ {i}) s _ {2} (\alpha_ {i} ^ {\prime}) = (\varepsilon \circ (s _ {1} \otimes s _ {2})) \biggl (\sum_ {i = 1} ^ {r} \alpha \alpha_ {i} \otimes \alpha_ {i} ^ {\prime} \biggr) \\ = (\varepsilon \circ (s _ {1} \otimes s _ {2})) \biggl (\sum_ {i = 1} ^ {r} \alpha_ {i} \otimes \alpha_ {i} ^ {\prime} \alpha \biggr) = \sum_ {i = 1} ^ {r} s _ {1} (\alpha_ {i}) s _ {2} (\alpha_ {i} ^ {\prime} \alpha) = x s _ {2} (\alpha). \end{array}
\]

The equality \(x^{-1}S_{1}x = S_{2}\) follows, giving assertion a).

Under the assumptions of b), suppose furthermore that \(r^{2}=0\) . In this case, the (A, A)-bimodule r is annihilated by r, and we therefore view it as an \((\overline{A},\overline{A})\) -bimodule. Choose a K-linear section \(\sigma\) of \(\pi\) . We have

\[
\alpha x = \sigma (\alpha) x \quad \text { and } \quad x \alpha = x \sigma (\alpha)\tag{19}
\]

for \(\alpha \in \overline{\mathrm{A}}\) and \(x\in \mathfrak{r}\) . Set

\[
\varphi (\alpha , \beta) = \sigma (\alpha \beta) - \sigma (\alpha) \sigma (\beta)\tag{20}
\]

for \(\alpha, \beta \in \overline{\mathbf{A}}\) . We have the relation \(\pi(\varphi(\alpha, \beta)) = \alpha\beta - \alpha\beta = 0\) for \(\alpha, \beta \in \overline{\mathbf{A}}\) . Therefore, \(\varphi\) defines an element of \(\mathrm{C}^2(\overline{\mathbf{A}}, \mathfrak{r})\) . Let \(\alpha, \beta, \gamma\) be elements of \(\overline{\mathbf{A}}\) ; in view of (19), we have

\[
\begin{array}{r l} \partial^ {2} \varphi (\alpha , \beta , \gamma) & = \alpha \varphi (\beta , \gamma) - \varphi (\alpha \beta , \gamma) + \varphi (\alpha , \beta \gamma) - \varphi (\alpha , \beta) \gamma \\ & = \sigma (\alpha) \varphi (\beta , \gamma) - \varphi (\alpha \beta , \gamma) + \varphi (\alpha , \beta \gamma) - \varphi (\alpha , \beta) \sigma (\gamma) \\ & = \sigma (\alpha) (\sigma (\beta \gamma) - \sigma (\beta) \sigma (\gamma)) - \sigma (\alpha \beta \gamma) + \sigma (\alpha \beta) \sigma (\gamma) + \sigma (\alpha \beta \gamma) \\ & - \sigma (\alpha) \sigma (\beta \gamma) - (\sigma (\alpha \beta) - \sigma (\alpha) \sigma (\beta)) \sigma (\gamma) \\ & = 0. \end{array}
\]

By Theorem 3 of VIII, p.~242, the K-module \(\mathrm{H}^{2}(\overline{\mathrm{A}},\mathfrak{r})\) is reduced to zero. Hence there exists an element \(\psi\) of \(\mathrm{C}^{1}(\overline{\mathrm{A}},\mathfrak{r})\) such that \(\partial^{1}\psi = \varphi\) , in other words, such that we have

\[
\varphi (\alpha , \beta) = \alpha \psi (\beta) - \psi (\alpha \beta) + \psi (\alpha) \beta \quad \text { for } \alpha , \beta \text { in } \overline {{\mathrm{A}}}.\tag{21}
\]

We have \(\psi(\alpha)\psi(\beta)=0\) because \(r^{2}\) is zero; from (19) and (20), we then deduce

\[
(\sigma + \psi) (\alpha \beta) = (\sigma + \psi) (\alpha) (\sigma + \psi) (\beta),\tag{22}
\]

so that the K-linear section \(\sigma +\psi\) of \(\pi\) belongs to \(\Sigma\) . Its image is a subalgebra \(S\) of \(A\) such that \(A = S + \mathfrak{r}\) .

Let us now prove the existence of S in the general case. We reason by induction on the least integer \(p \geqslant 1\) such that \(r^{p} = 0\) ; the case p = 1 is trivial. Suppose \(p \geqslant 2\) , and set \(A' = A/r^{p-1}\) . The radical \(r'\) of \(A'\) is equal to \(r/r^{p-1}\) (Proposition 5 of VIII, p.~155), so satisfies \(r'^{p-1} = 0\) , and the algebra \(A'/r'\) is isomorphic to \(\overline{A} = A/r\) and is therefore absolutely semisimple. By the induction hypothesis, there exists a subalgebra \(S'\) of \(A'\) such that \(A' = S' \oplus r'\) . Then \(S'\) is of the form \(A''/r^{p-1}\) , where \(A''\) is a subalgebra of A containing \(r^{p-1}\) , and we have

\[
\mathrm{A} = \mathrm{A} ^ {\prime \prime} + \mathfrak {r}, \quad \mathfrak {r} ^ {p - 1} = \mathrm{A} ^ {\prime \prime} \cap \mathfrak {r}.\tag{23}
\]

The algebra \(A^{\prime\prime}/r^{p-1}\) is isomorphic to \(A^{\prime}/r^{\prime}\) ; we have \((\mathfrak{r}^{p-1})^{2}=0\) , so \(r^{p-1}\) is the radical of \(A^{\prime\prime}\) . By the case we just treated, there exists a subalgebra S of \(A^{\prime\prime}\) such that \(A^{\prime\prime}=S\oplus r^{p-1}\) ; we deduce the relation \(A=S\oplus r\) from (23).

COROLLARY 1 (Wedderburn's theorem). --- Let K be a commutative field, A a K-algebra, and r the radical of A. Suppose that the K-algebra A/r is absolutely semisimple.

\begin{enumerate}
\def\labelenumi{\alph{enumi})}
\item
  Let \(S_1\) and \(S_2\) be subalgebras of A satisfying \(A = S_1 \oplus \mathfrak{r} = S_2 \oplus \mathfrak{r}\) . There exists a special automorphism of A transforming \(S_1\) into \(S_2\) .
\item
  If \(\mathfrak{r}\) is nilpotent, then there exists a subalgebra \(S\) of \(A\) satisfying \(A = S \oplus \mathfrak{r}\) .
\end{enumerate}

This follows from Proposition 7 and Theorem 2 of VIII, p.~238.

COROLLARY 2. --- Let A be a commutative algebra of finite degree over a perfect field K, and let r be its radical. There exists a unique subalgebra S of A such that \(A = S \oplus r\) . Moreover, S is isomorphic to a product of finitely many extensions of K of finite degree.

The K-algebra A/r is semisimple (VIII, p.~173, Proposition 1) and has finite degree; it is absolutely semisimple because the field K is perfect (VIII, p.~232, Theorem 1). Since the ideal r is nilpotent and A is commutative, the existence and uniqueness of S then follow from Corollary 1. Since S is semisimple and commutative and has finite degree, the last assertion is a consequence of Proposition 3 of VIII, p.~137.

Remarks. --- 1) The assumption that A/r is absolutely semisimple is essential in Corollary 1 (VIII, p.~246, Exercise 4).

\begin{enumerate}
\def\labelenumi{\arabic{enumi})}
\setcounter{enumi}{1}
\tightlist
\item
  Suppose that A is an Artinian algebra over the field K. If A is commutative, then we can show (VIII, p.~180, Exercise 9) that A is isomorphic to a product of algebras \(A_{1} \times \cdots \times A_{n}\) such that \(\mathrm{A}_{i}/\Re(\mathrm{A}_{i})\) is a field for every i. Conversely, if A is not commutative, then A may not be isomorphic to a product of algebras \(A_{1} \times \cdots \times A_{n}\) such that \(\mathrm{A}_{i}/\Re(\mathrm{A}_{i})\) is a simple ring for every i (VIII, p.~247, Exercise 5).
\end{enumerate}

\BourbakiExercises{Exercises}

\begin{enumerate}
\def\labelenumi{\arabic{enumi})}
\tightlist
\item
  Let \(\mathrm{K}\) be a commutative field and \(\mathrm{A}\) a K-algebra. We call an A-module M absolutely without radical if for every extension \(\mathrm{L}\) of \(\mathrm{K}\) , the \(\mathrm{A}_{(\mathrm{L})}\) -module \(\mathrm{M}_{(\mathrm{L})}\) is without radical. We say that the algebra \(\mathrm{A}\) is absolutely without radical if the A-module \(\mathrm{A}_s\) is absolutely without radical.
\end{enumerate}

\begin{enumerate}
\def\labelenumi{\alph{enumi})}
\item
  A commutative algebra that is absolutely without radical is separable (V, §15, No.~2, p.~119, Definition 1).
\item
  Every submodule of a module that is absolutely without radical is absolutely without radical; a direct sum of modules that are absolutely without radical is absolutely without radical.
\item
  Let M be a finitely generated A-module; then M is absolutely without radical if and only if there exists a perfect field P that is an algebraic extension of K such that \(\mathfrak{R}_{\mathrm{A}_{(\mathrm{P})}}(\mathrm{M}_{(\mathrm{P})}) = 0\) .
\item
  A semisimple A-module M is absolutely without radical if and only if every simple module belonging to the support of M is.
\item
  Let S be a simple A-module and D its commutant. Prove that the following properties are equivalent:
\end{enumerate}

\begin{enumerate}
\def\labelenumi{(\roman{enumi})}
\item
  The A-module S is absolutely without radical.
\item
  The K-algebra \(\mathrm{D}\) is absolutely without radical.
\item
  The center of \(\mathrm{D}\) is a separable extension of \(\mathrm{K}\) .
\end{enumerate}

\(f\) ) A semisimple algebra is absolutely without radical if and only if the center of each of its simple components is a separable extension of K.

\(g\) ) Prove that an A-module that is absolutely without radical and finite-dimensional over \(\mathbf{K}\) is absolutely semisimple.

\begin{enumerate}
\def\labelenumi{\arabic{enumi})}
\setcounter{enumi}{1}
\tightlist
\item
  Let \(A_{1}\) and \(A_{2}\) be K-algebras, A the algebra \(A_{1} \otimes_{K} A_{2}\) , \(M_{i}\) an \(A_{i}\) -module (for i = 1, 2), and M the A-module \(M_{1} \otimes_{K} M_{2}\) .
\end{enumerate}

\begin{enumerate}
\def\labelenumi{\alph{enumi})}
\item
  Suppose that the \(A_{1}\) -module \(M_{1}\) is absolutely without radical. Prove the inclusion \(\mathfrak{R}_{\mathrm{A}}(\mathrm{M}) \subset \mathrm{M}_{1} \otimes_{\mathrm{K}} \mathfrak{R}_{\mathrm{A}_{2}}(\mathrm{M}_{2})\) .
\item
  If the modules \(M_{1}\) and \(M_{2}\) are absolutely without radical, then so is the A-module M.
\end{enumerate}

\begin{enumerate}
\def\labelenumi{\arabic{enumi})}
\setcounter{enumi}{2}
\item
  Let K be a commutative field, and let A and B be simple K-algebras. Suppose that A is of finite rank over its center Z, that Z is an algebraic extension of K, and that the K-algebra B is absolutely without radical (Exercise 1). Prove that the ring \(A \otimes_{K} B\) is absolutely flat (VIII, p.~171, Exercise 27; observe that every element of \(A \otimes_{K} B\) belongs to a subalgebra of the form \(A_{1} \otimes_{K} B\) , where \(A_{1}\) is a simple algebra of finite rank over K).
\item
  Let \(\mathrm{K}\) be a commutative field of characteristic \(p > 0\) , and let \(a\) be an element of \(\mathrm{K} - \mathrm{K}^p\) . Denote the (finite-dimensional) K-algebra \(\mathrm{K}[\mathrm{X}] / ((\mathrm{X}^p - a)^2)\) by A. Prove that the radical \(\mathfrak{r}\) of A has square zero, that the K-algebra \(\mathrm{A} / \mathfrak{r}\) is isomorphic to the radical extension \(\mathrm{L} = \mathrm{K}[\mathrm{X}] / (\mathrm{X}^{p} - a)\) of K, and that A does not contain any K-subalgebras isomorphic to L.
\item
  Let \(\mathrm{K}\) be a commutative field and \(\mathrm{L}\) a finite set. Denote the algebra constructed in Exercise 17 of VIII, p.~169 by B, the central element \(1 - \sum_{\lambda} c_{\lambda} + \sum_{\lambda \neq \mu} r_{\lambda \mu}\) of \(\mathrm{B}\) by \(\varepsilon\) , and the algebra \(\mathrm{B} / \mathrm{B}\varepsilon\) by A. The K-algebra A is finite-dimensional, its radical \(\mathfrak{r}\) has square zero, and \(\mathrm{A} / \mathfrak{r}\) is isomorphic to \(\mathrm{K}^{\mathrm{L}}\) . Prove that A is not isomorphic to a product of two nonzero algebras.
\end{enumerate}

¶ 6) Let K be a commutative field and E an extension of K. Prove that the following properties are equivalent:

\begin{enumerate}
\def\labelenumi{(\roman{enumi})}
\item
  For every separable extension L of K, the ring \(E_{(L)}\) is semisimple.
\item
  The field \(\mathrm{E}\) is a purely inseparable extension of a separable extension of \(\mathrm{K}\) of finite degree.
\end{enumerate}

(Assuming (i), use the corollary of V, §17, No.~2, p.~140 and Exercise 4 of V, §2, p.~146 to prove that E is an algebraic extension of K. Then, prove that the separable closure \(E_{s}\) of K in E has finite degree. To do this, note that in the opposite case, the ring \(\mathrm{E}_{(\mathrm{L})}\) , where L is a separable closure of K, would contain families of orthogonal idempotents \((e_{1},\ldots,e_{n})\) with n arbitrarily large. Assuming (ii), prove that if \(\mathrm{E}_{s(\mathrm{L})}\) is isomorphic to \(\prod_{j=1}^{r}C_{j}\) , then \(\mathrm{E}_{(\mathrm{L})}\) is isomorphic to \(\prod_{j=1}^{r}D_{j}\) , where \(D_{j}\) is a purely inseparable extension of \(C_{j}\) .)

¶ 7) Let E be a p-radical extension of a commutative field K of characteristic p \textgreater{} 0. Suppose that the ring \(E \otimes_{K} E\) is Noetherian.

\begin{enumerate}
\def\labelenumi{\alph{enumi})}
\item
  Prove that there exists an integer \(k\) such that \(\mathrm{E} \subset \mathrm{K}^{p^{-k}}\) (reason by contradiction: construct a sequence \((x_{n})\) in \(\mathrm{E}\) such that \((x_{n} \otimes 1 - 1 \otimes x_{n})^{p^{n-1}} \neq 0\) and \((x_{n} \otimes 1 - 1 \otimes x_{n})^{p^n} = 0\) ).
\item
  Let \(\mathrm{F} = \mathrm{K}(\mathrm{E}^{p})\) , and let \((x_{i})_{i \in \mathrm{I}}\) be a p-basis of E over F (V, §13, No.~1, p.~98). Set \(y_{i} = x_{i} \otimes 1\) and \(z_{i} = x_{i} \otimes 1 - 1 \otimes x_{i}\) . Let \(\Lambda\) be the subset of \(\mathbf{N}^{(I)}\) consisting of the families \((\alpha_{i})_{i \in \mathrm{I}}\) such that \(\alpha_{i} < p\) for every i. Prove that the elements \(\prod_{i \in I} y_{i}^{\alpha_{i}} z_{i}^{\beta_{i}}\) for \((\alpha_{i})\) , \((\beta_{i}) \in \Lambda\) form a basis of the F-vector space \(E \otimes_{F} E\) . Deduce that I is finite (observe that the ring \(E \otimes_{F} E\) is Noetherian).
\item
  Conclude that E has finite degree over K (use Exercise 1, b) of V, §13, p.~170).
\end{enumerate}

¶ 8) Let E be an algebraic extension of a commutative field K. Prove that the ring \(E \otimes_{K} E\) is Noetherian if and only if E has finite degree over K (reduce to the case of a purely inseparable extension using the method of Exercise 6, and then apply Exercise 7).

¶ 9) Let K be a commutative field and A a K-algebra. Suppose that for every Artinian K-algebra B, the ring \(A \otimes_{K} B\) is Artinian. Prove that A is finite-dimensional over K (reduce to the case when A is semisimple by observing that the \(A/\Re(A)\) -modules \(\Re(A)^{k}/\Re(A)^{k+1}\) have finite length; then use Exercises 6 and 7).

\begin{enumerate}
\def\labelenumi{\arabic{enumi})}
\setcounter{enumi}{9}
\tightlist
\item
  Let K be a commutative ring, A a K-algebra, and P an (A,A)-bimodule. An extension of A by P is a triple \((B,i,p)\) , where B is a K-algebra, \(p:B\to A\) a surjective homomorphism of K-algebras, and \(i:P\to B\) an injective K-linear homomorphism with image Ker p, satisfying \(i(p(b)x)=bi(x)\) and \(i(xp(b))=i(x)b\) for \(x\in P\) and \(b\in B\) . An isomorphism from an extension \((B,i,p)\) to an extension \((B',i',p')\) is an isomorphism of K-algebras \(u:B\to B'\) such that \(p'\circ u=p\) and \(u\circ i=i'\) .
\end{enumerate}

\begin{enumerate}
\def\labelenumi{\alph{enumi})}
\item
  Let \((B, i, p)\) be an extension of A by P; the image of i is a two-sided ideal of B with square zero. Conversely, if \(p : B \to A\) is a surjective homomorphism of K-algebras whose kernel P has square zero and i is the canonical injection of P into B, then the triple \((B, i, p)\) is an extension of A by P.
\item
  Let \(i_{0}: P \to A \oplus P\) and \(p_{0}: A \oplus P \to A\) be the canonical mappings. There exists a unique K-algebra structure on \(A \oplus P\) such that \((A \oplus P, i_{0}, p_{0})\) is an extension of A by P. Prove that the automorphism group of this extension can be identified with the group of K-derivations from A to P.
\item
  Let \((B, i, p)\) be an extension of A by P. Prove that the following conditions are equivalent:
\end{enumerate}

\begin{enumerate}
\def\labelenumi{(\roman{enumi})}
\item
  The extension \((\mathrm{B},i,p)\) is isomorphic to the extension \((\mathrm{A}\oplus \mathrm{P},i_0,p_0)\) .
\item
  There exists a homomorphism of K-algebras \(s: \mathrm{A} \to \mathrm{B}\) such that \(p \circ s = \operatorname{Id}_{\mathrm{A}}\) .
\item
  There exists a subalgebra \(S\) of \(B\) such that \(B = S \oplus i(P)\) .
\end{enumerate}

If these conditions are satisfied, then we say that the extension \((B,i,p)\) is trivial. d) Let \((B,i,p)\) and \((B',i',p')\) be extensions of A by P. Let C be the subalgebra of \(B \times B'\) consisting of the elements \((b,b')\) such that \(p(b) = p'(b')\) , let c be the two-sided ideal of C consisting of the elements \((i(x), -i'(x))\) for \(x \in P\) , and let \(B''\) be the algebra C/c and \(\pi : C \to B''\) the canonical homomorphism. Let \(i'' : P \to B''\) and \(p'' : B'' \to A\) be the homomorphisms defined by \(i''(x) = \pi(i(x),0)\) and \(p''(\pi(b,b')) = p(b)\) for \(x \in P\) , \(b \in B\) , \(b' \in B'\) . Prove that \((B'', i'', p'')\) is an extension of A by P; it is called the amalgamated sum of \((B,i,p)\) and \((B',i',p')\) . Prove that we thus define a commutative group law on the set \(\mathrm{Ex}_{\mathrm{K}}(\mathrm{A},\mathrm{P})\) of isomorphism classes of extensions of A by P, whose neutral element is the class of the trivial extension. Moreover, the action of K on \(\mathrm{Ex}_{\mathrm{K}}(\mathrm{A},\mathrm{P})\) given by \(\lambda(\mathrm{B},i,p) = (\mathrm{B},\lambda^{-1}i,p)\) for \(\lambda \in K\) , \(\lambda \neq 0\) defines a K-module structure on \(\mathrm{Ex}_{\mathrm{K}}(\mathrm{A},\mathrm{P})\) .

\begin{enumerate}
\def\labelenumi{\alph{enumi})}
\setcounter{enumi}{4}
\tightlist
\item
  Let \(u: P \to P'\) be a homomorphism of (A, A)-bimodules, and let \((B, i, p)\) be an extension of A by P. Let \(B'\) be the amalgamated sum \(B \oplus_{P} P'\) (II, §1, p.~381, Exercise 5); for \(b \in B\) , \(x' \in P'\) , denote the class of \((b, x')\) in \(B'\) by \([b, x']\) . Let \(i': P' \to B'\) , \(p': B' \to A\) , and \(v: B \to B'\) be the mappings define by
\end{enumerate}

\[
i ^ {\prime} (x ^ {\prime}) = [ 0, x ^ {\prime} ], \qquad p ^ {\prime} ([ b, x ^ {\prime} ]) = p (b), \qquad v (b) = [ b, 0 ].
\]

Prove that there exists a unique algebra structure on \(B'\) such that \((\mathrm{B}', i', p')\) is an extension of A by \(P'\) and v an algebra homomorphism. We say that \((\mathrm{B}', i', p')\) is the extension of A by \(P'\) deduced from \((\mathrm{B}, i, p)\) via u; every extension \((\mathrm{C}, j, q)\) of A by \(P'\) for which there exists an algebra homomorphism \(w: B \to C\) satisfying \(q \circ w = p\) and \(w \circ i = j \circ v\) is isomorphic to \((\mathrm{B}', i', p')\) .

\(f\) ) The mapping \(\mathrm{Hom}_{(\mathrm{A},\mathrm{A})}(\mathrm{P},\mathrm{P}')\times \mathrm{Ex}_\mathrm{K}(\mathrm{A},\mathrm{P})\to \mathrm{Ex}_\mathrm{K}(\mathrm{A},\mathrm{P}')\) that sends a pair \((u,(B,i,p))\) to the class of the extension deduced from (B, \(i,p)\) via \(u\) is K-bilinear. \(g\) ) Let \(f:\mathrm{A}'\to \mathrm{A}\) be a K-algebra homomorphism, and let (B, \(i,p)\) be an extension of A by P. Define likewise the structure of an algebra on the fibered product \(\mathrm{B}'' = \mathrm{B}\times_{\mathrm{A}}\mathrm{A}'\) and an extension \((\mathrm{B}'' , i'', p'')\) of \(\mathrm{A}'\) by P, said to be deduced from (B, \(i,p)\) via \(f\) .

\begin{enumerate}
\def\labelenumi{\alph{enumi})}
\setcounter{enumi}{7}
\tightlist
\item
  Suppose, moreover, that the algebra A is commutative. Prove that the classes of extensions \((B,i,p)\) of A by P such that the algebra B is commutative form a K-submodule of \(\operatorname{Ex}_{\mathrm{K}}(\mathrm{A},\mathrm{P})\) , which we denote by \(\operatorname{Ex}_{\mathrm{K}}^{c}(\mathrm{A},\mathrm{P})\) .
\end{enumerate}

\begin{enumerate}
\def\labelenumi{\arabic{enumi})}
\setcounter{enumi}{10}
\tightlist
\item
  \begin{enumerate}
  \def\labelenumii{\alph{enumii})}
  \tightlist
  \item
    Let f be an element of \(\mathrm{C}^{2}(\mathrm{A},\mathrm{P})\) (VIII, p.~239). Prove that we define the structure of an algebra on the K-module \(A \oplus P\) by setting \((a,x)(a',x') = (aa',ax' + xa' + f(a,a'))\) . The canonical injection \(i: P \to A \oplus P\) and the canonical surjection \(p: A \oplus P \to A\) are ring homomorphisms, and the triple \((\mathrm{A} \oplus \mathrm{P}, i, p)\) is an extension of A by P (observe that we have \(f(a,1) = af(1,1)\) and \(f(1,a) = f(1,1)a\) for every \(a \in A\) ).
  \end{enumerate}
\end{enumerate}

\begin{enumerate}
\def\labelenumi{\alph{enumi})}
\setcounter{enumi}{1}
\item
  Prove that the isomorphism class of this extension depends only on the class of f in \(\mathrm{H}^{2}(\mathrm{A},\mathrm{P})\) . So the construction in a) defines an isomorphism of K-modules from \(\mathrm{H}^{2}(\mathrm{A},\mathrm{P})\) to the K-submodule of \(\mathrm{Ex}_{\mathrm{K}}(\mathrm{A},\mathrm{P})\) consisting of the extensions of A by P that are trivial as extensions of K-modules.
\item
  Suppose that the ring A is commutative; denote by \(\mathrm{C}^{2}(\mathrm{A},\mathrm{P})^{s}\) the K-submodule of \(\mathrm{C}^{2}(\mathrm{A},\mathrm{P})\) consisting of the elements f such that \(f(x,y)=f(y,x)\) for all x,y in A and by \(\mathrm{H}^{2}(\mathrm{A},\mathrm{P})^{s}\) the image of \(\mathrm{C}^{2}(\mathrm{A},\mathrm{P})^{s}\) in \(\mathrm{H}^{2}(\mathrm{A},\mathrm{P})\) . Prove that the isomorphism in b) induces an isomorphism from \(\mathrm{H}^{2}(\mathrm{A},\mathrm{P})^{s}\) to the K-submodule of \(\mathrm{Ex}_{\mathrm{K}}^{c}(\mathrm{A},\mathrm{P})\) consisting of the commutative extensions of A by P that are trivial as extensions of K-modules.
\end{enumerate}

\begin{enumerate}
\def\labelenumi{\arabic{enumi})}
\setcounter{enumi}{11}
\tightlist
\item
  Let A be a commutative ring, J an ideal of A, B the ring A/J, and P a B-module. a) Prove that the exact sequence \(0 \rightarrow J/J^{2} \rightarrow A/J^{2} \rightarrow B \rightarrow 0\) defines an extension of the A-algebra B by \(J/J^{2}\) (Exercise 10).
\end{enumerate}

\begin{enumerate}
\def\labelenumi{\alph{enumi})}
\setcounter{enumi}{1}
\tightlist
\item
  Let P be a B-module. Prove that the homomorphism \(\operatorname{Hom}_{\mathrm{B}}(\mathrm{J}/\mathrm{J}^{2},\mathrm{P}) \to \operatorname{Ex}_{\mathrm{A}}^{c}(\mathrm{B},\mathrm{P})\) that sends u to the extension deduced from the extension in a) via u is an isomorphism.
\end{enumerate}

\section{Central Simple Algebras}

In this section, K is a commutative field.

\subsection{Central Simple Algebras}

DEFINITION 1. --- We say that an algebra A over the field K is central if the mapping \(\lambda\mapsto\lambda1\) is a bijection from K to the center of A.

A central algebra is not reduced to 0. For every integer \(n \geqslant 1\) , the matrix K-algebra \(\mathbf{M}_{n}(\mathrm{K})\) is central (VIII, p.~83, Corollary 2) and simple (VIII, p.~120, Theorem 1). More generally, let D be a central K-algebra of finite degree; then \(\mathbf{M}_{n}(\mathrm{D})\) is also central. Let A be a simple ring; its center Z is a field (VIII, p.~121, Corollary 1), and A is therefore a central simple algebra over Z. If the field K is algebraically closed, then a simple algebra of finite degree over K is central (VIII, p.~122, Corollary 3). The opposite algebra of a central simple algebra is central simple.

Remarks. --- 1) Let A and B be K-algebras. If \(A \otimes_{K} B\) is a central simple algebra, then so are A and B (III, §4, No.~4, p.~468, Corollary of Proposition 6 and VIII, p.~221, Proposition 6). Conversely, if A and B are central simple algebras and if one of them has finite degree over K, then \(A \otimes_{K} B\) is a central simple algebra (III, §4, No.~4, p.~468, Corollary of Proposition 6 and VIII, p.~222, Corollary 2).

\begin{enumerate}
\def\labelenumi{\arabic{enumi})}
\setcounter{enumi}{1}
\item
  Let A be a K-algebra and L an extension of the field K. If the L-algebra \(A_{(L)}\) is central simple, then the K-algebra A is central simple. Conversely, if one of the degrees {[}A : K{]} and {[}L : K{]} is finite and if A is a central simple K-algebra, then the L-algebra \(A_{(L)}\) is central simple. This follows from Corollary 2 of VIII, p.~222.
\item
  Let A and B be Morita-equivalent K-algebras. The algebra A is central simple if and only if B is (VIII, p.~105, Proposition 6; p.~111, Corollary; and p.~102, Corollary 1).
\item
  In particular, if A is a central simple K-algebra and \(n \geqslant 1\) , then \(\mathbf{M}_{n}(\mathrm{A})\) is a central simple K-algebra (VIII, p.~102, Example 1).
\end{enumerate}

THEOREM 1. --- Let A be a K-algebra of finite degree. The following properties are equivalent:

\begin{enumerate}
\def\labelenumi{(\roman{enumi})}
\item
  The algebra A is central simple.
\item
  The algebra \(\mathbf{A}\) is central and without radical.
\item
  The canonical homomorphism from the K-algebra \(A \otimes_{K} A^{\circ}\) to the K-algebra \(\operatorname{End}_{K}(A)\) that sends \(a \otimes a'\) to the K-linear mapping \(x \mapsto axa'\) from A to A is bijective.
\item
  There exist an extension L of the field K and an integer \(n \geqslant 1\) such that the L-algebras \(\mathrm{A}_{(\mathrm{L})}\) and \(\mathbf{M}_{n}(\mathrm{L})\) are isomorphic.
\item
  For every separable closure \(\mathrm{K}'\) of \(\mathbf{K}\) , there exists an integer \(n \geqslant 1\) such that the \(\mathrm{K}'\) -algebras \(\mathrm{A}_{(\mathrm{K}')}\) and \(\mathbf{M}_n(\mathrm{K}')\) are isomorphic.
\item
  There exist a Galois extension L of the field K of finite degree and an integer \(n \geqslant 1\) such that the L-algebras \(\mathrm{A}_{(\mathrm{L})}\) and \(\mathbf{M}_n(\mathrm{L})\) are isomorphic.
\item
  There exist a K-algebra of finite degree D that is a field with center K and an integer \(n \geqslant 1\) such that the algebra A is isomorphic to the algebra \(\mathbf{M}_{n}(\mathrm{D})\) .
\end{enumerate}

A ring is simple if and only if it is semisimple and its center is a field (VIII, p.~143, Corollary of Proposition 10). Since A is an algebra of finite degree over the field K, it is a left Artinian ring; it is therefore semisimple if and only if it is without radical (VIII, p.~154, Proposition 4). The equivalence of (i) and (ii) follows.

Set \(E = A \otimes_{K} A^{\circ}\) and \(F = \text{End}_{K}(A)\) ; denote by \(\varphi\) the canonical homomorphism from E to F defined by the relation \(\varphi(a \otimes a')(x) = axa'\) for \(x, a, a'\) in A. If the algebra A is central simple, then so are \(A^{\circ}\) and, therefore, E (Remark 1), so \(\varphi\) is injective. Now, we have \([E : K] = [A : K]^{2} = [F : K]\) , so \(\varphi\) is bijective. Conversely, suppose that \(\varphi\) is bijective; since the algebra F is central simple (because it is isomorphic to a matrix algebra \(\mathbf{M}_{m}(K)\) ), so is E, and therefore so is A (Remark 1). We have thus proved the equivalence of (i) and (iii).

By Remark 4, assertion (vii) implies assertion (i). The converse implication follows from Corollary 3 of VIII, p.~122 and Corollary 2 of VIII, p.~83.

It is clear that (vi) implies (iv), and (iv) implies (i) by Remark 2.

It remains to prove the implications (i)⇒(v)⇒(vi). Suppose that A is central simple, and let \(K'\) be a separable closure of K (V, §7, No.~8, p.~45). Then \(\mathrm{A}_{(\mathrm{K}')}\) is a central simple algebra of finite degree over \(K'\) (VIII, p.~251, Remark 2). By the corollary of VIII, p.~231, there consequently exist an integer \(n \geqslant 1\) and an isomorphism of \(K'\) -algebras from \(\mathrm{A}_{(\mathrm{K}')}\) to \(\mathbf{M}_{n}(\mathrm{K}')\) ; observe that the \(K'\) -algebras \(\mathbf{M}_{n}(\mathrm{K}')\) and \(\mathbf{M}_{n}(\mathrm{K})_{(\mathrm{K}')}\) are isomorphic. By Lemma 4 of VIII, p.~234, there exists a subextension L of \(K'\) , finitely generated over K, such that the L-algebras \(\mathrm{A}_{(\mathrm{L})}\) and \(\mathbf{M}_{n}(\mathrm{K})_{(\mathrm{L})}\) are isomorphic. Then L is separable and of finite degree over K, so contained in a subextension \(L'\) of \(K'\) that is Galois and of finite degree over K (V, §10, No.~1, p.~57, Proposition 2). The \(L'\) -algebras \(\mathrm{A}_{(\mathrm{L}')}\) and \(\mathbf{M}_{n}(\mathrm{L}')\) are then isomorphic.

COROLLARY 1. --- Let A be a central simple algebra of finite degree over a separably closed field K. There exists an integer \(n \geqslant 1\) such that A is isomorphic to the matrix algebra \(\mathbf{M}_{n}(\mathrm{K})\) .

Indeed, every Galois extension of \(\mathrm{K}\) is equal to \(\mathrm{K}\) ; it suffices to apply the equivalence of properties (i) and (v) of Theorem 1.

COROLLARY 2. --- Let A be a central simple algebra of finite degree over K (for example, a field of finite degree over K with center K). There exists an integer \(n \geqslant 1\) such that \([A : K] = n^{2}\) .

Let \(\mathrm{L}\) be an extension of \(\mathrm{K}\) and \(n\) a strictly positive integer such that the L-algebras \(\mathrm{A}_{(\mathrm{L})}\) and \(\mathbf{M}_n(\mathrm{L})\) are isomorphic. We have

\[
[ \mathrm{A}: \mathrm{K} ] = [ \mathrm{A} _ {(\mathrm{L})}: \mathrm{L} ] = [ \mathbf {M} _ {n} (\mathrm{L}): \mathrm{L} ] = n ^ {2}.
\]

In the notation of Corollary 2, the integer n is called the reduced degree of A.

Remark 5. --- Let A be a central simple algebra of finite degree over K whose reduced degree is a prime number \(\ell\) . Then either A is a field, or A is isomorphic to \(\mathbf{M}_{\ell}(\mathrm{K})\) . Indeed, A is isomorphic to an algebra of the form \(\mathbf{M}_{n}(\mathrm{D})\) , where D is a field with center K, and we have

\[
\ell^ {2} = [ \mathrm{A}: \mathrm{K} ] = n ^ {2} [ \mathrm{D}: \mathrm{K} ];
\]

if A is not a field, then \(n \neq 1\) , so \(n = \ell\) and D = K.

\subsection{Two Lemmas on Bimodules}

Let A and B be rings. For any homomorphism f from B to A, we denote by \(A^{f}\) the (B,A)-bimodule with underlying right A-module \(A_{d}\) and external law for its left B-module structure given by \((b,a)\mapsto f(b)a\) .

Lemma 1. --- Let f and g be homomorphisms from B to A. The following conditions are equivalent:

\begin{enumerate}
\def\labelenumi{(\roman{enumi})}
\item
  The (B,A)-bimodules \(A^{f}\) and \(A^{g}\) are isomorphic.
\item
  There exists an inner automorphism (I, §8, No.~4, p.~102, Example 2) \(\theta\) of A such that \(g = \theta \circ f\) .
\end{enumerate}

The automorphisms of the right A-module \(A_{d}\) are the mappings \(x \mapsto ax\) , where a is an invertible element of A. Such an automorphism is a B-linear mapping from \(A^{f}\) to \(A^{g}\) if and only if we have

\[
g (b) a x = a f (b) x
\]

for every x in A and every b in B. This relation is equivalent to \(g(b) = af(b)a^{-1}\) for every b in B, that is, to \(g = \theta \circ f\) , where \(\theta\) is the inner automorphism \(x \mapsto axa^{-1}\) of A.

Lemma 2. --- Suppose that B is a semisimple ring that is a finitely generated module over its center Z. Let M and N be (B, A)-bimodules. Suppose that they have finite length (which is, in particular, the case when they are right A-modules of finite length). If M and N are isomorphic as (Z, A)-bimodules, then they are isomorphic as (B, A)-bimodules.

\begin{enumerate}
\def\labelenumi{\Alph{enumi})}
\item
  First consider the case when B is the endomorphism ring of a vector space S of finite dimension d over a commutative field L. We then have Z = L; we view S as a (B, Z)-bimodule. The ring B is simple, S is a simple B-module, and Z is the commutant of S; every B-module is isotypical of type S (VIII, p.~122, Proposition 2 a)). Let \((V, \alpha)\) (resp. \((W, \beta)\) ) be a description of the B-module M (resp. N). The set V (resp. W) is endowed with a (Z, A)-bimodule structure such that \(\alpha\) (resp. \(\beta\) ) is an isomorphism of (B, A)-bimodules (VIII, p.~64, Remark 2). As (Z, A)-bimodules, M is isomorphic to \(V^{d}\) and N to \(W^{d}\) , and there exists an isomorphism from the set of (Z, A)-sub-bimodules of V, ordered by inclusion, to that of the (B, A)-sub-bimodules of M (loc. cit.). Hence V is a (Z, A)-bimodule of finite length, and so is W. Since the (Z, A)-bimodules \(V^{d}\) and \(W^{d}\) are isomorphic, the (Z, A)-bimodules V and W are isomorphic by Theorem 2, d) of VIII, p.~37 applied to the ring \(Z \otimes_{Z} A^{\circ}\) . Finally, the (B, A)-bimodules M and N are isomorphic.
\item
  Now consider the case when B is a simple ring that is finitely generated as a Z-module. Then Z is a field, and B is a central simple algebra of finite degree over the field Z. By Theorem 1 of VIII, p.~252, there exists an extension \(Z'\) of Z of finite degree over Z such that the \(Z'\) -algebra \(B' = B_{(Z')}\) is isomorphic to the endomorphism algebra of a finite-dimensional \(Z'\) -vector space. Set \(M' = M_{(Z')}\) and \(N' = N_{(Z')}\) . Then \(M'\) and \(N'\) are \((B', A)\) -bimodules of finite length; viewed as \((Z', A)\) -bimodules, \(M'\) and \(N'\) are isomorphic. By the case treated in A), \(M'\) and \(N'\) are isomorphic as \((B', A)\) -bimodules and a fortiori as \((B, A)\) -bimodules. Set \(r = [Z' : Z]\) . The \((B, A)\) -bimodule \(M' = Z' \otimes_{Z} M\) is isomorphic to \(M^r\) , and, likewise, the \((B, A)\) -bimodule \(N'\) is isomorphic to \(N^r\) . Since M and N are \((B, A)\) -bimodules of finite length, it follows from Theorem 2, d) of VIII, p.~37 that the \((B, A)\) -bimodules M and N are isomorphic.
\item
  Finally, consider the general case, when B is a semisimple ring that is finitely generated as a Z-module. Let S be the set of classes of simple B-modules; it is finite (VIII, p.~136, Proposition 1). For any \(\lambda \in S\) , denote by \(M_{\lambda}\) (resp. \(N_{\lambda}\) ) the isotypical component of type \(\lambda\) of the B-module M (resp. N); it is a (B, A)-sub-bimodule of M (resp. N) (Remark, VIII, p.~67). For \(\lambda \in S\) , denote the annihilator of the B-module \(\lambda\) by \(b_{\lambda}\) , and set \(B_{\lambda} = B / b_{\lambda}\) ; let \(Z_{\lambda}\) be the center of \(B_{\lambda}\) . For \(\lambda \in S\) , the \((B_{\lambda}, A)\) -bimodules \(M_{\lambda}\) and \(N_{\lambda}\) have finite length. We can then identify B with the product of the simple rings \(B_{\lambda}\) and Z with the product of the \(Z_{\lambda}\) (VIII, p.~141, Proposition 8). Moreover, we can identify M with \(\prod_{\lambda \in S} M_{\lambda}\) and N with \(\prod_{\lambda \in S} N_{\lambda}\) . By assumption, M and N are isomorphic as (Z, A)-bimodules; it follows that for \(\lambda \in S\) , \(M_{\lambda}\) and \(N_{\lambda}\) are isomorphic ( \(Z_{\lambda}, A\) )-bimodules. By the case treated in B), the \((B_{\lambda}, A)\) -bimodules \(M_{\lambda}\) and \(N_{\lambda}\) are isomorphic, and so the (B, A)-bimodules M and N are isomorphic.
\end{enumerate}

Remark. --- It follows from the proof of Lemma 2 that M and N are (Z, A)-bimodules of finite length. Consequently, if B and A are two semisimple rings that are finitely generated modules over their respective centers Z(B) and Z(A), then two (B, A)-bimodules of finite length that are isomorphic as (Z(B), Z(A))-bimodules are isomorphic.
